\chapter{Relación 1: Concurrencia, Intercalación y Lógica de Hoare}
\label{cap:relacion1}

\cuadrofuentes{
    \textbf{Fuentes de referencia cotejadas:}\\
    $\bullet$ \textbf{Relación Oficial de Problemas de Teoría del Profesor} (Manuel I. Capel Tuñón, Dpto. de Lenguajes y Sistemas Informáticos, Universidad de Granada, Curso 2026/2027).\\
    $\bullet$ \textbf{Temario y Relaciones Resueltas de Infomates} (LosDelDGIIM, Sección 5.1: 49 ejercicios de concurrencia y verificación formal).\\
    $\bullet$ \textbf{Relación de Ejercicios Resueltos de Ismael Sallami Moreno} (Doble Grado Informática + ADE, UGR).\\
    $\bullet$ \textbf{Materiales de InfoAdeTech Hub} y colecciones de exámenes de la UGR.
}

\begin{aviso}[Estructura Integral de la Relación 1]
Esta relación contiene los \textbf{36 problemas oficiales} del Tema 1 de Sistemas Concurrentes y Distribuidos de la UGR completamente resueltos, demostrados formalmente y explicados paso a paso desde los fundamentos. 

Están organizados en \textbf{7 bloques temáticos} respetando el orden secuencial estricto ($1 \to 36$) para permitir una localización inmediata de cualquier ejercicio de la hoja del profesor:
\begin{itemize}[leftmargin=*,noitemsep]
    \item \textbf{Bloque I (Ej. 1--3):} Concurrencia básica, intercalación, condiciones de carrera y búferes.
    \item \textbf{Bloque II (Ej. 4--10):} Grafos de precedencia, Bernstein, cobegin/coend, fork/join y sincronización con banderas/mallas.
    \item \textbf{Bloque III (Ej. 11--13):} Semántica axiomática de Hoare, axioma de asignación y reglas de consecuencia.
    \item \textbf{Bloque IV (Ej. 14--21):} Concurrencia en lógica de Hoare, teoremas de no interferencia de Owicki-Gries e invariantes.
    \item \textbf{Bloque V (Ej. 22--27):} Reglas de derivación (conjunción, intercambio sin aux, condicional, bucles y aliasing).
    \item \textbf{Bloque VI (Ej. 28--29):} Asignación a arrays/vectores y cálculo de la precondición más débil ($wp$).
    \item \textbf{Bloque VII (Ej. 30--36):} Bucles, invariantes inductivos, algoritmos de búsqueda, combinatorios y función de cota de terminación.
\end{itemize}
\end{aviso}

% ==============================================================================
% BLOQUE I: CONCURRENCIA BÁSICA, INTERCALACIÓN Y CONDICIONES DE CARRERA
% ==============================================================================
\section{Bloque I: Concurrencia Básica, Intercalación y Condiciones de Carrera}

% ------------------------------------------------------------------------------
% EJERCICIO 1
% ------------------------------------------------------------------------------
\begin{ejercicio}[1: Posibles valores de variable compartida tras bucles concurrentes]
Considerar el siguiente fragmento de programa para 2 procesos $P_1$ y $P_2$. Los dos procesos pueden ejecutarse a cualquier velocidad. ¿Cuáles son los posibles valores resultantes para la variable $x$? Suponer que $x$ debe ser cargada en un registro para incrementarse y que cada proceso usa un registro diferente para realizar el incremento.

\begin{lstlisting}[language=C++]
{ variables compartidas }
var x : integer := 0;

Process P1;                  Process P2;
var i: integer;              var j: integer;
begin                        begin
  for i:= 1 to 2 do begin      for j:= 1 to 2 do begin
    x:= x + 1;                   x:= x + 1;
  end;                         end;
end;                         end;
\end{lstlisting}
\end{ejercicio}

\begin{solucion}
Cada proceso ejecuta un bucle de 2 iteraciones. En cada iteración, la sentencia $x := x + 1$ no es atómica y se descompone a nivel de arquitectura en tres acciones atómicas elementales:
\begin{enumerate}
    \item \textbf{Lectura ($l_{ij}$):} Carga del valor actual de la memoria compartida $x$ en el registro privado $R_i$.
    \item \textbf{Incremento:} Operación aritmética interna en la ALU: $R_i \leftarrow R_i + 1$.
    \item \textbf{Escritura ($e_{ij}$):} Almacenamiento del valor de $R_i$ de vuelta en la variable compartida $x$.
\end{enumerate}
En total, $P_1$ ejecuta la secuencia $\{l_{11}, e_{11}, l_{12}, e_{12}\}$ y $P_2$ ejecuta $\{l_{21}, e_{21}, l_{22}, e_{22}\}$.

\textbf{1. Valor Máximo Posible ($x = 4$):}
Ocurre cuando no hay solapamiento destructivo entre lecturas y escrituras (ejecución puramente secuencial o intercalada sin solapamiento de lecturas):
\[
\text{Traza: } l_{11}, e_{11} (x=1) \to l_{12}, e_{12} (x=2) \to l_{21}, e_{21} (x=3) \to l_{22}, e_{22} (x=4).
\]

\textbf{2. Valor Mínimo Posible ($x = 2$):}
Ocurre cuando las lecturas de ambos procesos se solapan en ambas iteraciones, anulando mutuamente sus incrementos:
\begin{center}
\small
\begin{tabular}{cccc}
\toprule
\textbf{Acción $P_1$} & \textbf{Acción $P_2$} & \textbf{Registros ($R_1, R_2$)} & \textbf{Valor de $x$} \\ \midrule
$l_{11}$ (lee 0) & $l_{21}$ (lee 0) & $R_1=0, R_2=0$ & 0 \\
$e_{11}$ (escribe 1) & --- & $R_1=1$ & 1 \\
--- & $e_{21}$ (escribe 1) & $R_2=1$ & 1 \\
$l_{12}$ (lee 1) & $l_{22}$ (lee 1) & $R_1=1, R_2=1$ & 1 \\
$e_{12}$ (escribe 2) & --- & $R_1=2$ & 2 \\
--- & $e_{22}$ (escribe 2) & $R_2=2$ & \textbf{2} \\ \bottomrule
\end{tabular}
\end{center}

\textbf{3. Valor Intermedio ($x = 3$):}
Ocurre si el solapamiento destructivo solo tiene lugar en una de las dos iteraciones mientras que la otra se ejecuta limpiamente de forma secuencial. Por ejemplo: $P_1$ y $P_2$ solapan la primera iteración alcanzando $x = 1$, y luego $P_1$ y $P_2$ ejecutan secuencialmente sus segundas iteraciones incrementando sucesivamente a 2 y 3.

\textbf{¿Pueden obtenerse valores menores que 2 (por ejemplo, $x = 0$ o $x = 1$)?}
No. Para que el programa finalice, ambos procesos deben terminar obligatoriamente sus dos iteraciones. Como $x_0 = 0$, cada proceso lee un valor no negativo, suma 1 y escribe; y ningún proceso decrementa $x$. Además, al menos el último proceso en escribir en su segunda iteración habrá leído un valor $\ge 1$ (obtenido tras el primer incremento), por lo que forzosamente escribirá un valor final $\ge 2$.

\textbf{Conclusión:} El conjunto exacto de valores posibles finales es $\mathbf{\{2, 3, 4\}}$.
\end{solucion}

% ------------------------------------------------------------------------------
% EJERCICIO 2
% ------------------------------------------------------------------------------
\begin{ejercicio}[2: Condiciones de carrera, determinismo y atomicidad]
Para cada uno de los siguientes programas concurrentes:
\begin{itemize}[leftmargin=*]
    \item Indicar si existe una condición de carrera (\textit{race condition}).
    \item Determinar si el resultado final es determinista. En caso contrario, obtener los posibles valores finales de las variables.
    \item Identificar qué operaciones de los procesos pueden interferir entre sí.
    \item Indicar, cuando sea necesario, qué instrucciones deberían ejecutarse de forma atómica para eliminar la condición de carrera. Razonar si ello es suficiente para obtener un resultado determinista.
\end{itemize}

\begin{minipage}{0.48\textwidth}
\textbf{Programa (a):}
\begin{lstlisting}[language=C++]
var x : integer := 0;
cobegin
    x := x + 1;
 ||
    x := x + 1;
coend;
\end{lstlisting}
\end{minipage}\hfill
\begin{minipage}{0.48\textwidth}
\textbf{Programa (b):}
\begin{lstlisting}[language=C++]
var x : integer := 0;
var y : integer := 0;
cobegin
    x := x + 1;
 ||
    y := y + 1;
coend;
\end{lstlisting}
\end{minipage}

\vspace{10pt}
\begin{minipage}{0.48\textwidth}
\textbf{Programa (c):}
\begin{lstlisting}[language=C++]
var x : integer := 1;
var y : integer := 0;
cobegin
    x := y + 1;
 ||
    y := x + 1;
coend;
\end{lstlisting}
\end{minipage}\hfill
\begin{minipage}{0.48\textwidth}
\textbf{Programa (d):}
\begin{lstlisting}[language=C++]
var x : integer := 2;
var y : integer := 3;
cobegin
    x := x * y;
 ||
    y := x * y;
coend;
\end{lstlisting}
\end{minipage}
\end{ejercicio}

\begin{solucion}
Analizamos cada programa sistemáticamente aplicando las condiciones de Bernstein y el modelo de intercalación de instrucciones a nivel de lectura/escritura en registros:

\subsection*{Programa (a)}
\begin{itemize}[leftmargin=*]
    \item \textbf{Condición de carrera:} \textbf{SÍ}. Ambos procesos acceden concurrentemente a la variable compartida $x$ y ambos realizan operaciones de escritura ($E(P_1) \cap E(P_2) = \{x\} \neq \emptyset$).
    \item \textbf{Determinismo:} \textbf{NO es determinista}. Si se solapan las lecturas ($l_1, l_2, e_1, e_2$), ambos leen 0 y escriben 1, resultando en $x = 1$. Si se ejecutan secuencialmente, el resultado es $x = 2$. Valores posibles: $\mathbf{x \in \{1, 2\}}$.
    \item \textbf{Operaciones que interfieren:} La lectura de un proceso con la escritura del otro ($l_1$ con $e_2$ y $l_2$ con $e_1$), así como las dos escrituras concurrentes.
    \item \textbf{Atomicidad:} Hacer atómicas ambas asignaciones ($\langle x := x + 1 \rangle$) es \textbf{suficiente para garantizar el determinismo}, ya que la suma es asociativa y conmutativa. Las dos únicas trazas atómicas posibles dan idéntico resultado: $x = 2$.
\end{itemize}

\subsection*{Programa (b)}
\begin{itemize}[leftmargin=*]
    \item \textbf{Condición de carrera:} \textbf{NO}. $P_1$ únicamente lee y escribe en $x$ ($L_1 = E_1 = \{x\}$), mientras que $P_2$ únicamente lee y escribe en $y$ ($L_2 = E_2 = \{y\}$). Los conjuntos de variables son completamente disjuntos:
    \[
    L_1 \cap E_2 = \emptyset, \quad E_1 \cap L_2 = \emptyset, \quad E_1 \cap E_2 = \emptyset.
    \]
    Se satisfacen estrictamente las tres condiciones de Bernstein.
    \item \textbf{Determinismo:} \textbf{SÍ es determinista}. El resultado final es siempre único e invariable: $\mathbf{x = 1, y = 1}$.
    \item \textbf{Interferencias:} Ninguna.
    \item \textbf{Atomicidad:} No es necesaria ninguna acción atómica adicional.
\end{itemize}

\subsection*{Programa (c)}
\begin{itemize}[leftmargin=*]
    \item \textbf{Condición de carrera:} \textbf{SÍ}. $P_1$ lee $y$ y escribe en $x$ ($L_1 = \{y\}, E_1 = \{x\}$). $P_2$ lee $x$ y escribe en $y$ ($L_2 = \{x\}, E_2 = \{y\}$). Se violan dos condiciones de Bernstein: $E_1 \cap L_2 = \{x\} \neq \emptyset$ y $L_1 \cap E_2 = \{y\} \neq \emptyset$.
    \item \textbf{Determinismo:} \textbf{NO es determinista}. Posibles valores finales:
    \begin{enumerate}
        \item Si $P_1$ termina antes de que $P_2$ lea: $x = 0 + 1 = 1$, luego $y = 1 + 1 = 2 \implies \mathbf{(x=1, y=2)}$.
        \item Si $P_2$ termina antes de que $P_1$ lea: $y = 1 + 1 = 2$, luego $x = 2 + 1 = 3 \implies \mathbf{(x=3, y=2)}$.
        \item Si ambos leen antes de que ninguno escriba: $P_1$ lee $y=0 \Rightarrow x=1$; $P_2$ lee $x=1 \Rightarrow y=2 \implies \mathbf{(x=1, y=2)}$.
    \end{enumerate}
    Conjunto de estados finales posibles: $\mathbf{\{(1, 2), (3, 2)\}}$.
    \item \textbf{Operaciones que interfieren:} La lectura de $x$ por $P_2$ interfiere con la escritura de $x$ por $P_1$; la lectura de $y$ por $P_1$ interfiere con la escritura de $y$ por $P_2$.
    \item \textbf{Atomicidad:} Hacer atómicas las instrucciones ($\langle x := y + 1 \rangle$ e $\langle y := x + 1 \rangle$) \textbf{NO es suficiente para lograr determinismo}, ya que subsiste una carrera de grano grueso entre qué instrucción atómica se ejecuta primero: si se ejecuta $P_1$ primero da $(1, 2)$, y si se ejecuta $P_2$ primero da $(3, 2)$. Para asegurar determinismo se requeriría imponer un orden de sincronización estricto.
\end{itemize}

\subsection*{Programa (d)}
\begin{itemize}[leftmargin=*]
    \item \textbf{Condición de carrera:} \textbf{SÍ}. Ambos procesos leen y escriben sobre las mismas variables compartidas ($L_1 = \{x, y\}, E_1 = \{x\}$; $L_2 = \{x, y\}, E_2 = \{y\}$).
    \item \textbf{Determinismo:} \textbf{NO es determinista}.
    \begin{enumerate}
        \item Si $P_1$ precede totalmente a $P_2$: $x = 2 \times 3 = 6$, luego $y = 6 \times 3 = 18 \implies \mathbf{(x=6, y=18)}$.
        \item Si $P_2$ precede totalmente a $P_1$: $y = 2 \times 3 = 6$, luego $x = 2 \times 6 = 12 \implies \mathbf{(x=12, y=6)}$.
        \item Si ambos leen inicialmente $x=2, y=3$ antes de escribir: $P_1$ calcula $2 \times 3 = 6$ y escribe $x=6$; $P_2$ calcula $2 \times 3 = 6$ y escribe $y=6 \implies \mathbf{(x=6, y=6)}$.
    \end{enumerate}
    Conjunto de estados finales posibles: $\mathbf{\{(6, 18), (12, 6), (6, 6)\}}$.
    \item \textbf{Atomicidad:} Hacer atómicas las operaciones elimina el estado $(6, 6)$, pero \textbf{NO elimina el no-determinismo}, pues persisten los resultados $(6, 18)$ y $(12, 6)$ según el orden de intercalación.
\end{itemize}
\end{solucion}

% ------------------------------------------------------------------------------
% EJERCICIO 3
% ------------------------------------------------------------------------------
\begin{ejercicio}[3: Copia concurrente de un fichero f en g con doble búfer]
¿Cómo se podría hacer la copia del fichero $f$ en otro $g$, de forma concurrente, usando dos búferes intermedios $B_1$ y $B_2$, de tal manera que el proceso que lee de $f$ y el proceso que escribe en $g$ puedan trabajar en paralelo?
\end{ejercicio}

\begin{solucion}
La técnica de \textbf{doble búfer} (\textit{double buffering}) permite solapar la operación de lectura de entrada/salida (I/O) desde el disco con la operación de procesamiento/escritura en destino:
\begin{itemize}
    \item Mientras el proceso \textbf{Lector} carga el siguiente bloque de datos desde el fichero $f$ en el búfer $B_1$, el proceso \textbf{Escritor} vuelca el contenido del bloque previo desde el búfer $B_2$ hacia el fichero $g$.
    \item En la siguiente iteración, ambos procesos intercambian los roles de los búferes de manera alternada y sincronizada.
\end{itemize}

\textbf{Implementación formal con banderas de sincronización / semáforos binarios:}
\begin{lstlisting}[language=C++]
// Recursos compartidos
var B1, B2 : TipoBloque;
var lleno1, lleno2 : boolean := false, false;
var fin_fichero : boolean := false;

Process Lector;
begin
  repeat
    // Fase 1: Leer sobre B1 mientras Escritor procesa B2
    if not fin_fichero then begin
      leer_bloque(f, B1, fin_fichero);
      lleno1 := true;
      esperar_a_que(lleno2 == false); // Esperar que B2 haya sido vaciado
    end;
    // Fase 2: Leer sobre B2 mientras Escritor procesa B1
    if not fin_fichero then begin
      leer_bloque(f, B2, fin_fichero);
      lleno2 := true;
      esperar_a_que(lleno1 == false); // Esperar que B1 haya sido vaciado
    end;
  until fin_fichero;
end;

Process Escritor;
begin
  repeat
    // Esperar y vaciar B1
    esperar_a_que(lleno1 == true or fin_fichero);
    if lleno1 then begin
      escribir_bloque(g, B1);
      lleno1 := false;
    end;
    // Esperar y vaciar B2
    esperar_a_que(lleno2 == true or fin_fichero);
    if lleno2 then begin
      escribir_bloque(g, B2);
      lleno2 := false;
    end;
  until fin_fichero and not lleno1 and not lleno2;
end;
\end{lstlisting}

\textbf{Ventaja clave:} El tiempo total de copia se reduce drásticamente de $T_{\text{sec}} = N \cdot (T_{\text{leer}} + T_{\text{escribir}})$ a $T_{\text{conc}} \approx N \cdot \max(T_{\text{leer}}, T_{\text{escribir}})$, maximizando el uso del bus y los canales DMA de entrada/salida.
\end{solucion}

% ==============================================================================
% BLOQUE II: GRAFOS DE PRECEDENCIA, SINCRONIZACIÓN Y PARALELISMO
% ==============================================================================
\section{Bloque II: Grafos de Precedencia, Sincronización y Paralelismo}

% ------------------------------------------------------------------------------
% EJERCICIO 4
% ------------------------------------------------------------------------------
\begin{ejercicio}[4: Construcción de programas con cobegin-coend y fork-join para grafos]
Construir, utilizando las instrucciones concurrentes \texttt{cobegin-coend} y \texttt{fork-join}, programas concurrentes que se correspondan con los dos grafos de precedencia siguientes:

\begin{center}
\begin{tikzpicture}[>=Stealth, node distance=1.5cm, auto, thick,
    task/.style={circle, draw=scdPrimary, fill=scdLightAccent, text=scdDarkBlue, font=\bfseries\sffamily, minimum size=8mm}]
    % Grafo G1
    \node (tit1) at (1.5, 2.5) {\textbf{Grafo $G_1$}};
    \node[task] (A) at (1.5, 1.5) {$S_1$};
    \node[task] (B) at (0.3, 0) {$S_2$};
    \node[task] (C) at (2.7, 0) {$S_3$};
    \node[task] (D) at (0.3, -1.5) {$S_4$};
    \node[task] (E) at (2.7, -1.5) {$S_5$};
    \node[task] (F) at (1.5, -3) {$S_6$};
    
    \draw[->] (A) -- (B);
    \draw[->] (A) -- (C);
    \draw[->] (B) -- (D);
    \draw[->] (C) -- (E);
    \draw[->] (D) -- (F);
    \draw[->] (E) -- (F);

    % Grafo G2
    \node (tit2) at (7, 2.5) {\textbf{Grafo $G_2$}};
    \node[task] (S1) at (7, 1.5) {$S_1$};
    \node[task] (S2) at (5.2, 0) {$S_2$};
    \node[task] (S3) at (7, 0) {$S_3$};
    \node[task] (S4) at (8.8, 0) {$S_4$};
    \node[task] (S5) at (6.1, -1.5) {$S_5$};
    \node[task] (S6) at (7.9, -1.5) {$S_6$};
    \node[task] (S7) at (7, -3) {$S_7$};

    \draw[->] (S1) -- (S2);
    \draw[->] (S1) -- (S3);
    \draw[->] (S1) -- (S4);
    \draw[->] (S2) -- (S5);
    \draw[->] (S3) -- (S5);
    \draw[->] (S3) -- (S6);
    \draw[->] (S4) -- (S6);
    \draw[->] (S5) -- (S7);
    \draw[->] (S6) -- (S7);
\end{tikzpicture}
\end{center}
\end{ejercicio}

\begin{solucion}
\subsection*{Análisis del Grafo $G_1$}
$G_1$ es un grafo bien anidado (serie-paralelo): $S_1$ precede a dos ramas paralelas disjuntas ($S_2 \to S_4$ y $S_3 \to S_5$), las cuales deben converger antes de que inicie $S_6$.
\begin{itemize}[leftmargin=*]
    \item \textbf{Con \texttt{cobegin-coend}:}
\begin{lstlisting}[language=C++]
S1;
cobegin
    begin S2; S4; end;
 ||
    begin S3; S5; end;
coend;
S6;
\end{lstlisting}

    \item \textbf{Con \texttt{fork-join}:}
\begin{lstlisting}[language=C++]
    var c : integer := 2;
    S1;
    fork L_rama2;
    // Rama 1
    S2;
    S4;
    goto L_convergencia;
L_rama2:
    S3;
    S5;
L_convergencia:
    join c;
    S6;
\end{lstlisting}
\end{itemize}

\subsection*{Análisis del Grafo $G_2$}
En $G_2$, $S_3$ bifurca su flujo hacia $S_5$ y $S_6$, pero $S_5$ depende de $\{S_2, S_3\}$ mientras que $S_6$ depende de $\{S_3, S_4\}$. Esto genera un cruce de dependencias interno no estrictamente anidado.
\begin{itemize}[leftmargin=*]
    \item \textbf{Con \texttt{fork-join} (expresión directa de grafos arbitrarios):}
\begin{lstlisting}[language=C++]
    var c5 : integer := 2; // S5 espera a S2 y S3
    var c6 : integer := 2; // S6 espera a S3 y S4
    var c7 : integer := 2; // S7 espera a S5 y S6

    S1;
    fork L_S2;
    fork L_S3;
    // Tarea S4 ejecutada por el hilo principal
    S4;
    join c6; // Si c6==0, pasa a ejecutar S6
    goto L_S6;

L_S2:
    S2;
    join c5;
    goto L_S5;

L_S3:
    S3;
    // S3 senala a ambos contadores
    // Si al decrementar c5 llega a 0, activa S5
    if (decrementar(c5) == 0) fork L_S5;
    if (decrementar(c6) == 0) fork L_S6;
    quit;

L_S5:
    S5;
    join c7;
    goto L_S7;

L_S6:
    S6;
    join c7;
    goto L_S7;

L_S7:
    S7;
\end{lstlisting}

    \item \textbf{Con \texttt{cobegin-coend} (usando variables de sincronización / semáforos):} Como no es estrictamente serie-paralelo, se requiere sincronización auxiliar (o duplicación estructurada):
\begin{lstlisting}[language=C++]
var finS3 : boolean := false;
S1;
cobegin
    begin
        S2;
        esperar_a_que(finS3);
        S5;
    end
 ||
    begin
        S3;
        finS3 := true;
    end
 ||
    begin
        S4;
        esperar_a_que(finS3);
        S6;
    end
coend;
S7;
\end{lstlisting}
\end{itemize}
\end{solucion}

% ------------------------------------------------------------------------------
% EJERCICIO 5
% ------------------------------------------------------------------------------
\begin{ejercicio}[5: Obtención de grafos de precedencia a partir de código]
Dados los siguientes fragmentos de programas concurrentes, obtener sus grafos de precedencia asociados:

\begin{minipage}{0.48\textwidth}
\textbf{Fragmento (a):}
\begin{lstlisting}[language=C++]
S1;
cobegin
    S2;
 ||
    begin
        S3;
        cobegin S4; || S5; coend;
        S6;
    end
 ||
    S7;
coend;
S8;
\end{lstlisting}
\end{minipage}\hfill
\begin{minipage}{0.48\textwidth}
\textbf{Fragmento (b):}
\begin{lstlisting}[language=C++]
    var c1, c2 : integer := 2;
    S1;
    fork L1;
    S2;
    goto L2;
L1: S3;
L2: join c1;
    fork L3;
    S4;
    goto L4;
L3: S5;
L4: join c2;
    S6;
\end{lstlisting}
\end{minipage}
\end{ejercicio}

\begin{solucion}
Analizamos el orden de ejecución y sincronizaciones de cada fragmento:

\subsection*{Grafo del Fragmento (a)}
\begin{enumerate}
    \item $S_1$ se ejecuta secuencialmente al principio y precede a las tres ramas del \texttt{cobegin}.
    \item El \texttt{cobegin} exterior lanza 3 ramas concurrentes:
    \begin{itemize}
        \item Rama 1: $S_2$.
        \item Rama 2: $S_3 \to$ (\texttt{cobegin} interno de $S_4 \parallel S_5$) $\to S_6$.
        \item Rama 3: $S_7$.
    \end{itemize}
    \item El \texttt{coend} exterior obliga a que finalicen $S_2$, $S_6$ y $S_7$ antes de ejecutar $S_8$.
\end{enumerate}

\begin{center}
\begin{tikzpicture}[>=Stealth, node distance=1.2cm, auto, thick,
    task/.style={circle, draw=scdPrimary, fill=scdLightAccent, text=scdDarkBlue, font=\bfseries\sffamily, minimum size=7mm}]
    \node[task] (S1) at (0, 3) {$S_1$};
    \node[task] (S2) at (-2.5, 1) {$S_2$};
    \node[task] (S3) at (0, 1.8) {$S_3$};
    \node[task] (S4) at (-0.8, 0.5) {$S_4$};
    \node[task] (S5) at (0.8, 0.5) {$S_5$};
    \node[task] (S6) at (0, -0.8) {$S_6$};
    \node[task] (S7) at (2.5, 1) {$S_7$};
    \node[task] (S8) at (0, -2.2) {$S_8$};

    \draw[->] (S1) -- (S2);
    \draw[->] (S1) -- (S3);
    \draw[->] (S1) -- (S7);
    \draw[->] (S3) -- (S4);
    \draw[->] (S3) -- (S5);
    \draw[->] (S4) -- (S6);
    \draw[->] (S5) -- (S6);
    \draw[->] (S2) -- (S8);
    \draw[->] (S6) -- (S8);
    \draw[->] (S7) -- (S8);
\end{tikzpicture}
\end{center}

\subsection*{Grafo del Fragmento (b)}
\begin{enumerate}
    \item $S_1$ inicia y realiza \texttt{fork L1}, creando dos hilos que ejecutan $S_2$ y $S_3$ en paralelo.
    \item En \texttt{L2}, la instrucción \texttt{join c1} (con contador inicial 2) sincroniza la finalización de ambos ($S_2$ y $S_3$).
    \item Una vez superado el primer join, se lanza \texttt{fork L3}, ejecutando en paralelo $S_4$ y $S_5$.
    \item En \texttt{L4}, la instrucción \texttt{join c2} sincroniza a $S_4$ y $S_5$.
    \item Finalmente se ejecuta secuencialmente $S_6$.
\end{enumerate}

\begin{center}
\begin{tikzpicture}[>=Stealth, node distance=1.2cm, auto, thick,
    task/.style={circle, draw=scdPrimary, fill=scdLightAccent, text=scdDarkBlue, font=\bfseries\sffamily, minimum size=7mm}]
    \node[task] (S1) at (0, 2.5) {$S_1$};
    \node[task] (S2) at (-1.2, 1.2) {$S_2$};
    \node[task] (S3) at (1.2, 1.2) {$S_3$};
    \node[task] (S4) at (-1.2, -0.3) {$S_4$};
    \node[task] (S5) at (1.2, -0.3) {$S_5$};
    \node[task] (S6) at (0, -1.6) {$S_6$};

    \draw[->] (S1) -- (S2);
    \draw[->] (S1) -- (S3);
    \draw[->] (S2) -- (S4);
    \draw[->] (S2) -- (S5);
    \draw[->] (S3) -- (S4);
    \draw[->] (S3) -- (S5);
    \draw[->] (S4) -- (S6);
    \draw[->] (S5) -- (S6);
\end{tikzpicture}
\end{center}
\end{solucion}

% ------------------------------------------------------------------------------
% EJERCICIO 6
% ------------------------------------------------------------------------------
\begin{ejercicio}[6: Sistema de tiempo real con captador de impulsos]
Suponer un sistema de tiempo real que dispone de un captador de impulsos conectado a un contador de energía eléctrica. La función del sistema consiste en contar el número de impulsos producidos en 1 hora (cada KWh consumido se cuenta como un impulso) e imprimir este número en un dispositivo de salida. Para ello se dispone de un programa concurrente con 2 procesos: un proceso acumulador (lleva la cuenta de los impulsos recibidos) y un proceso escritor (escribe en la impresora). En la variable común a los 2 procesos $n$ se lleva la cuenta de los impulsos. El proceso acumulador puede invocar un procedimiento \texttt{Espera\_impulso} para esperar a que llegue un impulso, y el proceso escritor puede llamar a \texttt{Espera\_fin\_hora} para esperar a que termine una hora. El código de los procesos de este programa podría ser el siguiente:

\begin{lstlisting}[language=C++]
var n : integer := 0;

Process Acumulador;          Process Escritor;
begin                        begin
  while true do begin          while true do begin
    Espera_impulso;              Espera_fin_hora;
    n := n + 1;                  imprimir(n);
  end;                           n := 0;
end;                           end;
                             end;
\end{lstlisting}
¿Existe algún problema de condición de carrera en este programa? Justificar la respuesta y proponer una solución correcta.
\end{ejercicio}

\begin{solucion}
\textbf{1. Identificación de Condiciones de Carrera:}
\textbf{SÍ, existen graves condiciones de carrera} debido al acceso no sincronizado a la variable compartida $n$:
\begin{itemize}[leftmargin=*]
    \item \textbf{Pérdida de impulsos durante el reseteo:} Cuando el proceso \texttt{Escritor} despierta tras una hora, lee $n$ para imprimirlo y a continuación ejecuta $n := 0$. Si un impulso físico llega justo entre la lectura para imprimir y la asignación $n := 0$, el proceso \texttt{Acumulador} ejecutará $n := n + 1$. Inmediatamente después, el \texttt{Escritor} sobreescribirá $n := 0$, \textbf{perdiendo para siempre el impulso recibido}.
    \item \textbf{No atomicidad en $n := n + 1$ frente a $n := 0$:} A nivel de registros de máquina, si la lectura de $n$ por el acumulador se solapa con la escritura de 0 por el escritor, el acumulador puede restaurar un valor obsoleto incrementado, falseando por completo la cuenta de la siguiente hora.
\end{itemize}

\textbf{2. Solución Correcta:}
Para resolverlo, la lectura del valor de la hora anterior y la puesta a cero de $n$ deben realizarse de forma \textbf{atómica} respecto al incremento, utilizando una variable local auxiliar y exclusión mutua:

\begin{lstlisting}[language=C++]
var n : integer := 0;
var cerrojo : Mutex;

Process Acumulador;
begin
  while true do begin
    Espera_impulso;
    cerrojo.lock();
    n := n + 1;
    cerrojo.unlock();
  end;
end;

Process Escritor;
var n_hora_anterior : integer;
begin
  while true do begin
    Espera_fin_hora;
    // Seccion critica: capturar y resetear de forma indivisible
    cerrojo.lock();
    n_hora_anterior := n;
    n := 0;
    cerrojo.unlock();
    
    imprimir(n_hora_anterior);
  end;
end;
\end{lstlisting}
De este modo, si llega un impulso durante la impresión, el acumulador esperará a que termine la sección crítica y lo sumará sobre el nuevo contador en cero, garantizando conservación estricta de todos los impulsos.
\end{solucion}

% ------------------------------------------------------------------------------
% EJERCICIO 7
% ------------------------------------------------------------------------------
\begin{ejercicio}[7: Producto y suma concurrente de vectores]
Supongamos un programa concurrente en el cual hay, en memoria compartida dos vectores $a$ y $b$ de enteros y con tamaño par, declarados como sigue:
\begin{lstlisting}[language=C++]
var a, b : array[1..2*n] of integer; { n es una constante predefinida }
\end{lstlisting}
Queremos escribir un programa para obtener en $b$ una copia ordenada del contenido de $a$. Para ello disponemos de la función \texttt{Sort} que ordena un tramo de $a$ (entre las entradas $s$ y $t$, ambas incluidas) y deja el resultado en el mismo tramo de $a$, y del procedimiento \texttt{Merge} que mezcla de forma ordenada dos tramos ordenados consecutivos de $a$ (desde $s$ hasta $m$, y desde $m+1$ hasta $t$) y escribe el resultado en el tramo correspondiente de $b$ (desde $s$ hasta $t$). Escribir un programa concurrente que ordene el vector $a$ en paralelo aprovechando al máximo el multiprocesamiento.
\end{ejercicio}

\begin{solucion}
Dado que el vector tiene tamaño par $2n$, el enfoque óptimo de división y conquista (\textit{Divide and Conquer}) consiste en:
\begin{enumerate}
    \item Dividir el vector $a$ en dos mitades disjuntas: el tramo inferior $a[1..n]$ y el tramo superior $a[n+1..2n]$.
    \item Lanzar dos procesos concurrentes con \texttt{cobegin-coend}, ordenando cada mitad en paralelo mediante la función \texttt{Sort}.
    \item Una vez completadas ambas ordenaciones (garantizado por el \texttt{coend}), mezclar ambas mitades ordenadas en el vector $b[1..2n]$ usando \texttt{Merge}.
\end{enumerate}

\textbf{Código del programa concurrente:}
\begin{lstlisting}[language=C++]
program OrdenacionParalela;
const n = 1000;
var a, b : array[1..2*n] of integer;

begin
  // Fase 1: Ordenacion en paralelo de las dos mitades disjuntas
  cobegin
    Sort(1, n);       // Proceso 1 ordena a[1..n]
 ||
    Sort(n+1, 2*n);   // Proceso 2 ordena a[n+1..2*n]
  coend;

  // Fase 2: Mezcla ordenada secuencial de ambos tramos en b
  Merge(1, n, 2*n);   // Mezcla tramos [1..n] y [n+1..2*n] dejando el resultado en b[1..2*n]
end.
\end{lstlisting}

\textbf{Comprobación de Condiciones de Bernstein en la Fase Paralela:}
\begin{itemize}
    \item Proceso 1 ($P_1$): Lee y escribe únicamente en el rango de memoria $a[1..n]$. Por tanto: $L(P_1) = E(P_1) = \{a[1], \dots, a[n]\}$.
    \item Proceso 2 ($P_2$): Lee y escribe únicamente en el rango $a[n+1..2n]$. Por tanto: $L(P_2) = E(P_2) = \{a[n+1], \dots, a[2n]\}$.
    \item Intersecciones:
    \[
    L(P_1) \cap E(P_2) = \emptyset, \quad E(P_1) \cap L(P_2) = \emptyset, \quad E(P_1) \cap E(P_2) = \emptyset.
    \]
\end{itemize}
Se satisfacen estrictamente las tres condiciones de Bernstein, garantizando ausencia de interferencias y determinismo absoluto.
\end{solucion}

% ------------------------------------------------------------------------------
% EJERCICIO 8
% ------------------------------------------------------------------------------
\begin{ejercicio}[8: Multiplicación de matrices concurrentes y condiciones de Bernstein]
Supongamos que tenemos un programa con tres matrices ($a, b$ y $c$) de valores flotantes declaradas como variables globales. La multiplicación secuencial de $a$ y $b$ (almacenando el resultado en $c$) se puede hacer mediante un procedimiento \texttt{MultiplicacionSec} declarado como aparece aquí:

\begin{lstlisting}[language=C++]
var a, b, c : array[1..3, 1..3] of real;

procedure MultiplicacionSec()
var i, j, k : integer;
begin
  for i := 1 to 3 do
    for j := 1 to 3 do begin
      c[i,j] := 0;
      for k := 1 to 3 do
        c[i,j] := c[i,j] + a[i,k] * b[k,j];
    end;
end;
\end{lstlisting}
Escribir un programa con el mismo fin, pero que use 3 procesos concurrentes. Suponer que los elementos de las matrices $a$ y $b$ se pueden leer simultáneamente, así como que elementos distintos de $c$ pueden escribirse simultáneamente.
\end{ejercicio}

\begin{solucion}
\textbf{1. Análisis de Bernstein por filas:}
Para realizar la multiplicación con 3 procesos concurrentes ($P_1, P_2, P_3$), asignamos a cada proceso el cálculo completo de una fila entera de la matriz resultado $c$:
\begin{itemize}
    \item Proceso $P_i$ calcula la fila $i$ ($c[i, 1], c[i, 2], c[i, 3]$).
    \item \textbf{Conjunto de lectura de $P_i$:} Para calcular la fila $i$, $P_i$ necesita los elementos de la fila $i$ de la matriz $a$ y todas las columnas de la matriz $b$:
    \[
    L(P_i) = \{a[i, 1], a[i, 2], a[i, 3]\} \cup \{b[k, j] \mid 1 \le k, j \le 3\}.
    \]
    \item \textbf{Conjunto de escritura de $P_i$:} $P_i$ escribe única y exclusivamente en los elementos de la fila $i$ de $c$:
    \[
    E(P_i) = \{c[i, 1], c[i, 2], c[i, 3]\}.
    \]
\end{itemize}

\textbf{2. Comprobación de las Condiciones de Bernstein ($i \neq j$):}
\begin{enumerate}
    \item $E(P_i) \cap E(P_j) = \{c[i, *]\} \cap \{c[j, *]\} = \emptyset$, pues pertenecen a filas distintas de $c$.
    \item $L(P_i) \cap E(P_j) = (a \cup b) \cap \{c[j, *]\} = \emptyset$, pues $c$ no es leída por ningún proceso.
    \item $E(P_i) \cap L(P_j) = \{c[i, *]\} \cap (a \cup b) = \emptyset$.
\end{enumerate}
Dado que las lecturas simultáneas sobre $a$ y $b$ están permitidas por el hardware y no generan interferencia ($L(P_i) \cap L(P_j) \neq \emptyset$ es inocuo), los 3 procesos son totalmente independientes.

\textbf{3. Código Concurrente con 3 Procesos:}
\begin{lstlisting}[language=C++]
procedure MultiplicarFila(i : integer);
var j, k : integer;
begin
  for j := 1 to 3 do begin
    c[i,j] := 0;
    for k := 1 to 3 do
      c[i,j] := c[i,j] + a[i,k] * b[k,j];
  end;
end;

procedure MultiplicacionConc();
begin
  cobegin
    MultiplicarFila(1); // Proceso P1 calcula fila 1
 ||
    MultiplicarFila(2); // Proceso P2 calcula fila 2
 ||
    MultiplicarFila(3); // Proceso P3 calcula fila 3
  coend;
end;
\end{lstlisting}
\textit{Nota de paralelismo máximo:} Se podrían utilizar hasta 9 procesos concurrentes asignando un proceso independiente a cada celda individual $c[i,j]$, ya que todos los $E(P_{ij}) = \{c[i,j]\}$ son disjuntos dos a dos.
\end{solucion}

% ------------------------------------------------------------------------------
% EJERCICIO 9
% ------------------------------------------------------------------------------
\begin{ejercicio}[9: Grafo de dependencias de 9 actividades con EsperarPor y Acabar]
Un trozo de programa ejecuta nueve procesos o actividades ($P_1, P_2, \dots, P_9$), repetidas veces, de forma concurrente con \texttt{cobegin-coend}, pero que requieren sincronizarse según un grafo de dependencias en malla $3 \times 3$:

\begin{center}
\begin{tikzpicture}[>=Stealth, node distance=1.8cm, auto, thick,
    task/.style={circle, draw=scdPrimary, fill=scdLightAccent, text=scdDarkBlue, font=\bfseries\sffamily, minimum size=9mm}]
    % Nodos 3x3
    \node[task] (P1) at (0, 3) {$P_1$};
    \node[task] (P2) at (2.2, 3) {$P_2$};
    \node[task] (P3) at (4.4, 3) {$P_3$};

    \node[task] (P4) at (0, 1.5) {$P_4$};
    \node[task] (P5) at (2.2, 1.5) {$P_5$};
    \node[task] (P6) at (4.4, 1.5) {$P_6$};

    \node[task] (P7) at (0, 0) {$P_7$};
    \node[task] (P8) at (2.2, 0) {$P_8$};
    \node[task] (P9) at (4.4, 0) {$P_9$};

    % Dependencias horizontales
    \draw[->] (P1) -- (P2);
    \draw[->] (P2) -- (P3);
    \draw[->] (P4) -- (P5);
    \draw[->] (P5) -- (P6);
    \draw[->] (P7) -- (P8);
    \draw[->] (P8) -- (P9);

    % Dependencias verticales
    \draw[->] (P1) -- (P4);
    \draw[->] (P4) -- (P7);
    \draw[->] (P2) -- (P5);
    \draw[->] (P5) -- (P8);
    \draw[->] (P3) -- (P6);
    \draw[->] (P6) -- (P9);
\end{tikzpicture}
\end{center}

\begin{lstlisting}[language=C++]
while true do
  cobegin
    P1; P2; P3; P4; P5; P6; P7; P8; P9;
  coend;
\end{lstlisting}
Supón que queremos realizar la sincronización indicada en el grafo usando llamadas desde cada rutina a dos procedimientos:
\begin{itemize}[leftmargin=*]
    \item \texttt{EsperarPor(i)}: llamado por una rutina cualquiera (la número $k$) para esperar a que termine la rutina número $i$, usando espera ocupada.
    \item \texttt{Acabar(i)}: llamado por la rutina número $i$, al final de la misma, para indicar que dicha rutina ya ha finalizado.
\end{itemize}
Escribe una implementación de \texttt{EsperarPor} y \texttt{Acabar} (junto con la declaración e inicialización de las variables compartidas necesarias) que cumpla con los requisitos dados.
\end{ejercicio}

\begin{solucion}
\textbf{1. Diseño del vector de estados compartidos:}
Para señalizar la terminación de cada proceso se utiliza un vector global booleano \texttt{finalizado[1..9]}.
\begin{itemize}
    \item \texttt{finalizado[i] == true} indica que el proceso $P_i$ ha completado su ejecución en la iteración actual.
    \item \texttt{finalizado[i] == false} indica que el proceso $P_i$ aún no ha terminado.
\end{itemize}

\textbf{2. Reto clave: Reinicialización cíclica segura:}
Dado que el programa está envuelto en un bucle infinito \texttt{while true do cobegin ... coend}, el vector \texttt{finalizado} debe restablecerse a \texttt{false} para la siguiente iteración. Si no se reinicia, en la iteración 2 todos los \texttt{EsperarPor} evaluarían a verdadero inmediatamente, perdiéndose la sincronización.
Dado que $P_9$ es el último nodo del grafo (sumidero hacia el cual convergen directa o indirectamente todos los flujos), cuando $P_9$ finaliza se garantiza que \textbf{todos los procesos $P_1, \dots, P_8$ ya han finalizado}. Por tanto, $P_9$ es el único responsable de reiniciar el vector a \texttt{false}.

\textbf{3. Implementación de los Procedimientos:}
\begin{lstlisting}[language=C++]
// Variable compartida en memoria comun
var finalizado : array[1..9] of boolean := (false, false, false, false, false, false, false, false, false);

procedure EsperarPor(i : integer);
begin
  // Espera ocupada (polling) hasta que el proceso predecesor senale finalizacion
  while not finalizado[i] do begin
    // No-op (espera activa)
  end;
end;

procedure Acabar(i : integer);
var k : integer;
begin
  if i < 9 then begin
    finalizado[i] := true; // Senala que Pi ha concluido
  end else begin
    // El proceso P9 es el sumidero final de la malla.
    // Antes de terminar el cobegin, reinicia el vector para la siguiente vuelta del while.
    for k := 1 to 9 do
      finalizado[k] := false;
  end;
end;
\end{lstlisting}
\end{solucion}

% ------------------------------------------------------------------------------
% EJERCICIO 10
% ------------------------------------------------------------------------------
\begin{ejercicio}[10: Vector de procesos estáticos P[n: 1..9] para la malla de actividades]
En el problema anterior los procesos $P_1, P_2, \dots, P_9$ se ponían en marcha usando \texttt{cobegin/coend}. Escribe un programa equivalente que use declaración estática de procesos, mediante un vector de procesos $P$, con índices desde 1 hasta 9, ambos incluidos. El proceso $P[n]$ contiene una secuencia de instrucciones desconocida, que llamamos $S_n$, y además debe incluir las llamadas necesarias a \texttt{Acabar} y \texttt{EsperarPor} (con la misma implementación que antes) para lograr la sincronización adecuada. Se incluye aquí una plantilla:

\begin{lstlisting}[language=C++]
Process P[ n : 1..9 ]
begin
  ..... { esperar (si es necesario) a los procesos que corresponda }
  Sn ;  { sentencias especificas de este proceso (desconocidas) }
  ..... { senalar que hemos terminado }
end;
\end{lstlisting}
\end{ejercicio}

\begin{solucion}
En la topología en malla $3 \times 3$, cada nodo $P_n$ tiene como máximo dos predecesores directos de los que depende:
\begin{itemize}
    \item \textbf{Predecesor vertical (superior):} $P_{n-3}$. Solo existe si el nodo no está en la primera fila ($n > 3$).
    \item \textbf{Predecesor horizontal (izquierdo):} $P_{n-1}$. Solo existe si el nodo no está en la primera columna ($(n \pmod 3) \neq 1$).
\end{itemize}

\subsection*{Alternativa 1: Determinación de dependencias por Aritmética Modular (Óptima)}
\begin{lstlisting}[language=C++]
Process P[ n : 1..9 ]
begin
  while true do begin
    // 1. Espera al predecesor de arriba (si no esta en la primera fila)
    if (n > 3) then
      EsperarPor(n - 3);

    // 2. Espera al predecesor de la izquierda (si no esta en la primera columna)
    if (n mod 3) <> 1 then
      EsperarPor(n - 1);

    // 3. Ejecucion del trabajo especifico del proceso n
    Sn;

    // 4. Senalizacion de finalizacion
    Acabar(n);
  end;
end;
\end{lstlisting}

\subsection*{Alternativa 2: Matriz Genérica de Dependencias}
Para topologías arbitrarias que no sigan un patrón geométrico regular, se precalcula una matriz estática \texttt{espera[1..9, 1..2]} con los IDs de los predecesores (usando $-1$ para indicar ausencia de predecesor):
\begin{lstlisting}[language=C++]
var espera : array[1..9, 1..2] of integer := (
  (-1, -1),   // P1: no espera a nadie (fuente)
  (1, -1),    // P2: espera a P1
  (2, -1),    // P3: espera a P2
  (1, -1),    // P4: espera a P1
  (2, 4),     // P5: espera a P2 y P4
  (3, 5),     // P6: espera a P3 y P5
  (4, -1),    // P7: espera a P4
  (5, 7),     // P8: espera a P5 y P7
  (6, 8)      // P9: espera a P6 y P8 (sumidero)
);

Process P[ n : 1..9 ]
var i : integer;
begin
  while true do begin
    for i := 1 to 2 do
      if espera[n, i] <> -1 then
        EsperarPor(espera[n, i]);

    Sn;
    Acabar(n);
  end;
end;
\end{lstlisting}
\end{solucion}

% ==============================================================================
% BLOQUE III: SEMÁNTICA AXIOMÁTICA SECUENCIAL Y REGLAS DE HOARE
% ==============================================================================
\section{Bloque III: Semántica Axiomática Secuencial y Reglas de Hoare}

% ------------------------------------------------------------------------------
% EJERCICIO 11
% ------------------------------------------------------------------------------
\begin{ejercicio}[11: Obtención de poscondiciones adecuadas en Hoare]
Para los siguientes fragmentos de código, obtener la poscondición adecuada para convertirlo en un triple demostrable con la Lógica de Programas de Hoare:
\begin{enumerate}[label=(\alph*),leftmargin=*]
    \item $\{i < 10\} \; i := 2 * i + 1; \; \{ \; \}$
    \item $\{i > 0\} \; i := i - 1; \; \{ \; \}$
    \item $\{i > j\} \; i := i + 1; \; j := j + 1; \; \{ \; \}$
    \item $\{\text{falso}\} \; a := a + 7; \; \{ \; \}$
    \item $\{\text{verdad}\} \; i := 3; \; j := 2 * i; \; \{ \; \}$
    \item $\{\text{verdad}\} \; c := a + b; \; c := c / 2; \; \{ \; \}$
\end{enumerate}
\end{ejercicio}

\begin{solucion}
Para deducir la poscondición más fuerte y natural en cada caso, aplicamos el \textbf{Axioma de Asignación}:
\[
\{P[e/x]\} \; x := e \; \{P\}
\]
junto con la regla de composición secuencial para programas de varias instrucciones:

\begin{enumerate}[label=(\alph*),leftmargin=*]
    \item \textbf{Fragmento (a):} Si $i_{\text{ant}} < 10$, tras la asignación $i_{\text{nuevo}} = 2 \cdot i_{\text{ant}} + 1$, despejando tenemos $i_{\text{ant}} = (i_{\text{nuevo}} - 1)/2$. Sustituyendo en la precondición:
    \[
    \frac{i - 1}{2} < 10 \iff i - 1 < 20 \iff i < 21.
    \]
    Por tanto, el triple demostrable es:
    \[
    \mathbf{\{i < 10\} \; i := 2 * i + 1 \; \{i < 21\}}.
    \]

    \item \textbf{Fragmento (b):} Si $i_{\text{ant}} > 0$ (con $i \in \mathbb{Z}$ significa $i_{\text{ant}} \ge 1$), tras $i := i - 1$:
    \[
    i_{\text{ant}} = i_{\text{nuevo}} + 1 > 0 \iff i_{\text{nuevo}} > -1 \iff i_{\text{nuevo}} \ge 0.
    \]
    Por tanto:
    \[
    \mathbf{\{i > 0\} \; i := i - 1 \; \{i \ge 0\}}.
    \]

    \item \textbf{Fragmento (c):} Aplicamos composición secuencial con aserto intermedio:
    \[
    \{i > j\} \; i := i + 1 \; \{i > j + 1\} \; j := j + 1 \; \{i > j\}.
    \]
    Dado que ambos se incrementan en 1, la relación de orden estricto se preserva:
    \[
    \mathbf{\{i > j\} \; i := i + 1; \; j := j + 1 \; \{i > j\}}.
    \]

    \item \textbf{Fragmento (d):} La precondición es $\{\text{falso}\}$. Por el principio lógico \textit{ex falso quodlibet} (de una contradicción se sigue cualquier cosa), el triple es demostrable para \textbf{cualquier poscondición} $Q$, en particular para $\{\text{verdad}\}$ o $\{\text{falso}\}$:
    \[
    \mathbf{\{\text{falso}\} \; a := a + 7 \; \{\text{verdad}\}}.
    \]

    \item \textbf{Fragmento (e):}
    \[
    \{\text{verdad}\} \; i := 3 \; \{i = 3\} \; j := 2 * i \; \{i = 3 \land j = 6\}.
    \]
    Poscondición: $\mathbf{\{i = 3 \land j = 6\}}$.

    \item \textbf{Fragmento (f):}
    \[
    \{\text{verdad}\} \; c := a + b \; \{c = a + b\} \; c := c / 2 \; \left\{c = \frac{a + b}{2}\right\}.
    \]
    Poscondición: $\mathbf{\{c = (a + b) / 2\}}$.
\end{enumerate}
\end{solucion}

% ------------------------------------------------------------------------------
% EJERCICIO 12
% ------------------------------------------------------------------------------
\begin{ejercicio}[12: Triples no demostrables con la Lógica de Programas]
¿Cuáles de los siguientes triples \textbf{no son demostrables} con la Lógica de Programas?
\begin{enumerate}[label=(\alph*),leftmargin=*]
    \item $\{i > 0\} \; i := i - 1; \; \{i \ge 0\}$
    \item $\{x \ge 7\} \; x := x + 3; \; \{x \ge 9\}$
    \item $\{i < 9\} \; i := 2 * i + 1; \; \{i \le 20\}$
    \item $\{a > 0\} \; a := a - 7; \; \{a > -6\}$
\end{enumerate}
\end{ejercicio}

\begin{solucion}
Para cada triple $\{P\} \; x := e \; \{Q\}$, calculamos la precondición más débil requerida por el axioma de asignación $Q[e/x]$, y comprobamos si la precondición dada $P$ implica lógicamente $Q[e/x]$ ($P \implies Q[e/x]$):

\begin{itemize}[leftmargin=*]
    \item \textbf{Triple (a):} $Q \equiv \{i \ge 0\}$.
    \[
    Q[i - 1 / i] \equiv \{i - 1 \ge 0\} \equiv \{i \ge 1\} \equiv \{i > 0\} \quad (\text{para } i \in \mathbb{Z}).
    \]
    Como $\{i > 0\} \implies \{i \ge 1\}$, el triple \textbf{SÍ es demostrable}.

    \item \textbf{Triple (b):} $Q \equiv \{x \ge 9\}$.
    \[
    Q[x + 3 / x] \equiv \{x + 3 \ge 9\} \equiv \{x \ge 6\}.
    \]
    Como la precondición dada es $\{x \ge 7\}$ y trivialmente $\{x \ge 7\} \implies \{x \ge 6\}$, por la regla de fortalecimiento de precondición el triple \textbf{SÍ es demostrable}.

    \item \textbf{Triple (c):} $Q \equiv \{i \le 20\}$.
    \[
    Q[2i + 1 / i] \equiv \{2i + 1 \le 20\} \iff \{2i \le 19\} \iff \{i \le 9.5\} \iff \{i \le 9\}.
    \]
    Como la precondición dada es $\{i < 9\} \implies \{i \le 8\} \implies \{i \le 9\}$, por la regla de fortalecimiento el triple \textbf{SÍ es demostrable}.

    \item \textbf{Triple (d):} $Q \equiv \{a > -6\}$.
    \[
    Q[a - 7 / a] \equiv \{a - 7 > -6\} \iff \{a > 1\}.
    \]
    La precondición dada es $\{a > 0\}$. Para que fuese demostrable, se requeriría $\{a > 0\} \implies \{a > 1\}$, lo cual es \textbf{FALSO}:
    \[
    \text{\textbf{Contraejemplo:}} \quad \text{Si } a = 1, \text{ se cumple la precondición } 1 > 0.
    \]
    Tras ejecutar $a := a - 7$, el valor final es $a = 1 - 7 = -6$. Al evaluar la poscondición: $-6 > -6$ es \textbf{falso}.
\end{itemize}

\textbf{Conclusión:} El único triple que \textbf{NO es demostrable} es el \textbf{(d)}.
\end{solucion}

% ------------------------------------------------------------------------------
% EJERCICIO 13
% ------------------------------------------------------------------------------
\begin{ejercicio}[13: Reglas de consecuencia y deducibilidad de triples]
Si el triple $\{P\} \; C \; \{Q\}$ es demostrable, ¿cuál de los siguientes triples \textbf{no se puede demostrar} en general?
\begin{enumerate}[label=(\alph*),leftmargin=*]
    \item $\{P \land D\} \; C \; \{Q\}$
    \item $\{P \lor D\} \; C \; \{Q\}$
    \item $\{P\} \; C \; \{Q \lor D\}$
    \item $\{P\} \; C \; \{Q \lor P\}$
\end{enumerate}
\end{ejercicio}

\begin{solucion}
Las dos \textbf{Reglas de Consecuencia} de la lógica de Hoare establecen:
\begin{enumerate}
    \item \textbf{Fortalecimiento de la Precondición:}
    \[
    \frac{P' \implies P, \quad \{P\} \; C \; \{Q\}}{\{P'\} \; C \; \{Q\}}
    \]
    \item \textbf{Debilitamiento de la Poscondición:}
    \[
    \frac{\{P\} \; C \; \{Q\}, \quad Q \implies Q'}{\{P\} \; C \; \{Q'\}}
    \]
\end{enumerate}

Analizamos cada opción:
\begin{itemize}[leftmargin=*]
    \item \textbf{Opción (a):} $\{P \land D\} \implies P$ es una tautología lógica. Por la regla de fortalecimiento de precondición, el triple $\{P \land D\} \; C \; \{Q\}$ \textbf{siempre se puede demostrar}.
    \item \textbf{Opción (b):} Requiere $(P \lor D) \implies P$, lo cual exige $D \implies P$. Si $D$ no implica $P$ (por ejemplo, estados donde $D$ es verdadero pero $P$ es falso), no tenemos ninguna garantía sobre el comportamiento de $C$. Por tanto, $\{P \lor D\} \; C \; \{Q\}$ \textbf{NO se puede demostrar en general}.
    \item \textbf{Opción (c):} $Q \implies (Q \lor D)$ es una tautología lógica. Por la regla de debilitamiento de poscondición, el triple \textbf{siempre se puede demostrar}.
    \item \textbf{Opción (d):} Análogamente, $Q \implies (Q \lor P)$ es una tautología lógica. Por debilitamiento de poscondición, el triple \textbf{siempre se puede demostrar}.
\end{itemize}

\textbf{Conclusión:} El triple que \textbf{no se puede demostrar} es el \textbf{(b)}.
\end{solucion}

% ==============================================================================
% BLOQUE IV: CONCURRENCIA EN LÓGICA DE HOARE Y OWICKI-GRIES
% ==============================================================================
\section{Bloque IV: Concurrencia en Lógica de Hoare y Teorema de Owicki-Gries}

% ------------------------------------------------------------------------------
% EJERCICIO 14
% ------------------------------------------------------------------------------
\begin{ejercicio}[14: Valores finales con acciones atómicas cobegin]
Dado el programa concurrente:
\begin{lstlisting}[language=C++]
int x = 5, y = 2;
cobegin
    < x = x + y >;
 ||
    < y = x * y >;
coend;
\end{lstlisting}
obtener:
\begin{enumerate}[label=(\alph*),leftmargin=*]
    \item Los valores finales de $x$ e $y$.
    \item Los valores finales de $x$ e $y$ si quitamos los símbolos de instrucción atómica $\langle \; \rangle$.
\end{enumerate}
\end{ejercicio}

\begin{solucion}
\subsection*{Apartado (a): Con instrucciones atómicas $\langle \; \rangle$}
Como las instrucciones son atómicas, solo existen $2! = 2$ intercalaciones posibles a nivel de sentencias:

\begin{enumerate}[leftmargin=*]
    \item \textbf{Orden 1 ($P_1 \to P_2$):}
    \begin{itemize}
        \item Se ejecuta $P_1$: $x = 5 + 2 = 7$.
        \item Se ejecuta $P_2$: $y = 7 \times 2 = 14$.
        \item Resultado final: $\mathbf{(x = 7, y = 14)}$.
    \end{itemize}
    \item \textbf{Orden 2 ($P_2 \to P_1$):}
    \begin{itemize}
        \item Se ejecuta $P_2$: $y = 5 \times 2 = 10$.
        \item Se ejecuta $P_1$: $x = 5 + 10 = 15$.
        \item Resultado final: $\mathbf{(x = 15, y = 10)}$.
    \end{itemize}
\end{enumerate}
Conjunto de estados finales con atomicidad: $\mathbf{\{(7, 14), (15, 10)\}}$.

\subsection*{Apartado (b): Sin instrucciones atómicas (descomposición en registros)}
Cada asignación se descompone en lectura de operandos en registros locales, cálculo y escritura en memoria:
\begin{align*}
P_1: \quad & l_1^x: R_1 \leftarrow x; \quad l_1^y: R_2 \leftarrow y; \quad e_1^x: x \leftarrow R_1 + R_2. \\
P_2: \quad & l_2^x: R_3 \leftarrow x; \quad l_2^y: R_4 \leftarrow y; \quad e_2^y: y \leftarrow R_3 \times R_4.
\end{align*}
Analizamos las posibles intercalaciones no atómicas:
\begin{itemize}[leftmargin=*]
    \item Además de los dos órdenes puramente secuenciales anteriores que dan $(7, 14)$ y $(15, 10)$:
    \item \textbf{Lecturas cruzadas simultáneas antes de cualquier escritura:}
    Ambos procesos leen los valores iniciales $x = 5, y = 2$:
    \begin{itemize}
        \item $P_1$ lee $x=5, y=2 \implies$ calcula $R_1 + R_2 = 7$.
        \item $P_2$ lee $x=5, y=2 \implies$ calcula $R_3 \times R_4 = 10$.
        \item Sin importar en qué orden escriban ($e_1^x$ y luego $e_2^y$, o viceversa):
        $x$ recibe 7 e $y$ recibe 10.
        \item Resultado final: $\mathbf{(x = 7, y = 10)}$.
    \end{itemize}
\end{itemize}
Conjunto completo de estados finales posibles sin atomicidad: $\mathbf{\{(7, 14), (15, 10), (7, 10)\}}$.
\end{solucion}

% ------------------------------------------------------------------------------
% EJERCICIO 15
% ------------------------------------------------------------------------------
\begin{ejercicio}[15: Comprobación formal de No Interferencia de Owicki-Gries]
Comprobar si la demostración del triple:
\[
\{x \ge 2\} \; \langle x := x - 2 \rangle \; \{x \ge 0\}
\]
interfiere con los siguientes teoremas auxiliares:
\begin{enumerate}[label=(\alph*),leftmargin=*]
    \item $\{x \ge 0\} \; \langle x := x + 3 \rangle \; \{x \ge 3\}$
    \item $\{x \ge 0\} \; \langle x := x + 3 \rangle \; \{x \ge 0\}$
    \item $\{x \ge 7\} \; \langle x := x + 3 \rangle \; \{x \ge 10\}$
    \item $\{y \ge 0\} \; \langle y := y + 3 \rangle \; \{y \ge 3\}$
\end{enumerate}
\end{ejercicio}

\begin{solucion}
\textbf{Definición del Criterio de No Interferencia de Owicki-Gries:}
Dado un triple $\{P\} \; S \; \{Q\}$ y otra instrucción concurrente $S'$ con asertos asociados pre($S'$) y post($S'$), decimos que $S'$ \textbf{no interfiere} con el triple $\{P\} \; S \; \{Q\}$ si las aserciones $P$ y $Q$ se preservan invariantes bajo la ejecución de $S'$, es decir, si son demostrables los triples de no interferencia:
\[
\text{NI}_1: \quad \{P \land \text{pre}(S')\} \; S' \; \{P\}
\qquad\text{y}\qquad
\text{NI}_2: \quad \{Q \land \text{pre}(S')\} \; S' \; \{Q\}.
\]

En nuestro problema: $P \equiv \{x \ge 2\}$ y $Q \equiv \{x \ge 0\}$. Debemos comprobar si cada acción candidata $S'$ preserva $P$ y $Q$:

\begin{itemize}[leftmargin=*]
    \item \textbf{Candidato (a):} $S' \equiv \langle x := x + 3 \rangle$, con precondición $\text{pre}(S') \equiv \{x \ge 0\}$.
    \begin{enumerate}
        \item Comprobamos preservación de $P$:
        \[
        \{x \ge 2 \land x \ge 0\} \; x := x + 3 \; \{x \ge 2\}.
        \]
        Por el axioma de asignación: $(x + 3 \ge 2) \iff x \ge -1$. Como $(x \ge 2 \land x \ge 0) \equiv (x \ge 2) \implies (x \ge -1)$, se cumple.
        \item Comprobamos preservación de $Q$:
        \[
        \{x \ge 0 \land x \ge 0\} \; x := x + 3 \; \{x \ge 0\}.
        \]
        Por el axioma: $(x + 3 \ge 0) \iff x \ge -3$. Como $(x \ge 0) \implies (x \ge -3)$, se cumple.
    \end{enumerate}
    Por tanto, el candidato (a) \textbf{NO INTERFIERE}.

    \item \textbf{Candidato (b):} Idéntico a (a) porque la instrucción es la misma y la precondición es también $\{x \ge 0\}$. \textbf{NO INTERFIERE}.

    \item \textbf{Candidato (c):} $S' \equiv \langle x := x + 3 \rangle$, con precondición $\{x \ge 7\}$.
    Como $\{x \ge 7\} \implies \{x \ge 0\}$, las precondiciones conjuntas son aún más fuertes, por lo que también \textbf{NO INTERFIERE}.

    \item \textbf{Candidato (d):} $S' \equiv \langle y := y + 3 \rangle$ opera sobre una variable totalmente disjunta $y$.
    Como las aserciones $P$ y $Q$ dependen solo de $x$ y $S'$ solo escribe en $y$, por el axioma de asignación con variables disjuntas la no interferencia es inmediata. \textbf{NO INTERFIERE}.
\end{itemize}

\textbf{Conclusión:} \textbf{Ninguno de los cuatro teoremas interfiere} con el triple dado, pues todos incrementan $x$ (o modifican otra variable $y$), favoreciendo el mantenimiento de las cotas inferiores $x \ge 2$ y $x \ge 0$.
\end{solucion}

% ------------------------------------------------------------------------------
% EJERCICIO 16
% ------------------------------------------------------------------------------
\begin{ejercicio}[16: Triple concurrente con tres sumandos atómicos]
Dado el siguiente triple:
\begin{lstlisting}[language=C++]
{ x == 0 }
cobegin
    < x = x + a >;
 ||
    < x = x + b >;
 ||
    < x = x + c >;
coend
{ x == a + b + c }
\end{lstlisting}
demostrar que es correcto utilizando la Lógica de Hoare Concurrente.
\end{ejercicio}

\begin{solucion}
Para demostrar la corrección del programa concurrente mediante el Teorema de Owicki-Gries, debemos asociar asertos a cada proceso y demostrar:
\begin{enumerate}
    \item La corrección local de cada proceso por separado.
    \item La no interferencia mutua entre las aserciones de cada proceso y las acciones atómicas de los demás.
\end{enumerate}

\textbf{1. Definición de Variables Auxiliares de Estado:}
Introducimos variables booleanas auxiliares conceptuales $done_a, done_b, done_c$, inicializadas a \texttt{false}, que indican qué sumandos se han incorporado ya a $x$.
El \textbf{Invariante Global} del sistema es:
\[
I \equiv \left( x = (done_a ? a : 0) + (done_b ? b : 0) + (done_c ? c : 0) \right).
\]

\textbf{2. Demostración Local de cada Proceso:}
Para el proceso $P_a$ (y análogamente para $P_b$ y $P_c$):
\[
\text{pre}(P_a) \equiv I \land \neg done_a, \qquad \text{post}(P_a) \equiv I \land done_a.
\]
La instrucción atómica es $\langle x := x + a; \; done_a := true \rangle$.
Aplicando el axioma de asignación compuesto:
\[
(I \land done_a)[x + a / x, true / done_a] \equiv I \land \neg done_a.
\]
Por tanto, el triple local $\{I \land \neg done_a\} \; P_a \; \{I \land done_a\}$ es estrictamente demostrable.

\textbf{3. Demostración de No Interferencia de Owicki-Gries:}
Debemos comprobar que la ejecución de $P_b$ no destruye las aserciones de $P_a$:
\[
\{(I \land \neg done_a) \land \text{pre}(P_b)\} \; \langle x := x + b; \; done_b := true \rangle \; \{I \land \neg done_a\}.
\]
Dado que $P_b$ únicamente modifica $x$ sumándole $b$ y conmuta $done_b$ a \texttt{true}, el estado resultante mantiene $\neg done_a$ inalterado y actualiza $I$ de forma consistente. La misma invariancia se cumple para la poscondición $\{I \land done_a\}$.

\textbf{4. Conclusión de la Poscondición Final:}
Al alcanzar el \texttt{coend}, se cumple la conjunción de las poscondiciones locales:
\[
Q_{\text{final}} \equiv (I \land done_a \land done_b \land done_c) \implies (x = a + b + c).
\]
Queda demostrada la corrección total del triple concurrente.
\end{solucion}

% ------------------------------------------------------------------------------
% EJERCICIO 17
% ------------------------------------------------------------------------------
\begin{ejercicio}[17: Aplicación de las reglas de consecuencia de Hoare]
Suponer que el triple $\{suma > 1\} \; suma = suma + 4 \; \{suma > 5\}$ es demostrable. ¿Cuál de los siguientes triples es también demostrable? Indicar por qué.
\begin{enumerate}[label=(\alph*),leftmargin=*]
    \item $\{suma > 2\} \; suma = suma + 4 \; \{suma > 5\}$
    \item $\{suma \ge 1\} \; suma = suma + 4 \; \{suma > 5\}$
    \item $\{suma > 0\} \; suma = suma + 4 \; \{suma > 5\}$
    \item $\{suma > 1\} \; suma = suma + 4 \; \{suma > 6\}$
\end{enumerate}
\end{ejercicio}

\begin{solucion}
Nos basamos en las dos reglas de consecuencia de Hoare:
\begin{align*}
\text{Regla 1 (Fortalecer Precondición):} & \quad \frac{P' \implies P, \quad \{P\} S \{Q\}}{\{P'\} S \{Q\}} \\[6pt]
\text{Regla 2 (Debilitar Poscondición):} & \quad \frac{\{P\} S \{Q\}, \quad Q \implies Q'}{\{P\} S \{Q'\}}
\end{align*}

Analizamos las cuatro opciones:
\begin{itemize}[leftmargin=*]
    \item \textbf{Opción (a):} La nueva precondición es $\{suma > 2\}$.
    Es un hecho matemático que $(suma > 2) \implies (suma > 1)$.
    Por la regla de fortalecimiento de precondición, el triple $\{suma > 2\} \; suma = suma + 4 \; \{suma > 5\}$ \textbf{SÍ es demostrable}.
    \item \textbf{Opción (b):} La nueva precondición es $\{suma \ge 1\}$. Sin embargo, $(suma \ge 1) \not\implies (suma > 1)$ (contraejemplo: $suma = 1 \implies 1 + 4 = 5 \not> 5$). Falso.
    \item \textbf{Opción (c):} La nueva precondición es $\{suma > 0\}$, que tampoco implica $suma > 1$ (si $suma = 1$, daría $5 \not> 5$). Falso.
    \item \textbf{Opción (d):} La nueva poscondición es $\{suma > 6\}$. Para aplicar debilitamiento de poscondición se necesitaría $Q \implies Q'$, es decir, $(suma > 5) \implies (suma > 6)$, lo cual es falso (por ejemplo, $suma = 5.5$ o $suma = 6$). Falso.
\end{itemize}

\textbf{Conclusión:} El único triple demostrable es el \textbf{(a)}.
\end{solucion}

% ------------------------------------------------------------------------------
% EJERCICIO 18
% ------------------------------------------------------------------------------
\begin{ejercicio}[18: Selección de valores finales tras tres asignaciones atómicas]
Seleccionar el valor correcto de las 2 variables ($x$ e $y$) después de ejecutarse el siguiente programa concurrente:
\begin{lstlisting}[language=C++]
int x = 5, y = 2;
cobegin
    < x = x + y >;
 ||
    < y = x * y >;
 ||
    < x = x - y >;
coend;
\end{lstlisting}
Opciones disponibles:
(a) $x = 7, y = 14$ \qquad (b) $x = 5, y = 10$ \qquad (c) $x = -7, y = 14$ \qquad (d) $x = -3, y = 10$.
\end{ejercicio}

\begin{solucion}
Denotamos las 3 instrucciones atómicas como:
\[
S_1 \equiv \langle x := x + y \rangle, \qquad S_2 \equiv \langle y := x * y \rangle, \qquad S_3 \equiv \langle x := x - y \rangle.
\]
Partiendo de $x_0 = 5, y_0 = 2$, evaluamos las $3! = 6$ permutaciones atómicas posibles:

\begin{center}
\footnotesize
\setlength{\tabcolsep}{4.5pt}
\begin{tabular}{clllc}
\toprule
\textbf{Permutación} & \textbf{Paso 1} & \textbf{Paso 2} & \textbf{Paso 3} & \textbf{Resultado $(x, y)$} \\ \midrule
$S_1 \to S_2 \to S_3$ & $x = 5+2 = 7$ & $y = 7 \times 2 = 14$ & $x = 7 - 14 = -7$ & $\mathbf{(-7, 14)}$ \quad [Opción c] \\
$S_1 \to S_3 \to S_2$ & $x = 5+2 = 7$ & $x = 7 - 2 = 5$ & $y = 5 \times 2 = 10$ & $\mathbf{(5, 10)}$ \quad [Opción b] \\
$S_3 \to S_1 \to S_2$ & $x = 5-2 = 3$ & $x = 3 + 2 = 5$ & $y = 5 \times 2 = 10$ & $\mathbf{(5, 10)}$ \quad [Opción b] \\
$S_2 \to S_1 \to S_3$ & $y = 5 \times 2 = 10$ & $x = 5 + 10 = 15$ & $x = 15 - 10 = 5$ & $\mathbf{(5, 10)}$ \quad [Opción b] \\
$S_2 \to S_3 \to S_1$ & $y = 5 \times 2 = 10$ & $x = 5 - 10 = -5$ & $x = -5 + 10 = 5$ & $\mathbf{(5, 10)}$ \quad [Opción b] \\
$S_3 \to S_2 \to S_1$ & $x = 5-2 = 3$ & $y = 3 \times 2 = 6$ & $x = 3 + 6 = 9$ & $\mathbf{(9, 6)}$ \\ \bottomrule
\end{tabular}
\end{center}

\textbf{Conclusión:} De las opciones presentadas en el enunciado oficial de tipo test, las respuestas correctas son \textbf{(b)} ($x = 5, y = 10$) y \textbf{(c)} ($x = -7, y = 14$).
\end{solucion}

% ------------------------------------------------------------------------------
% EJERCICIO 19
% ------------------------------------------------------------------------------
\begin{ejercicio}[19: Traza de asertos y valores finales en programa secuencial]
Estudiar cuáles son los valores finales de las variables $x$ e $y$ en el siguiente programa secuencial. Insertar los asertos adecuados entre llaves, antes y después de cada sentencia, para poder obtener una traza de demostración formal del programa, que incluya en su último aserto los valores finales de las variables.
\begin{lstlisting}[language=C++]
int x = c1;
int y = c2;
x = x + y;
y = x * y;
x = x - y;
\end{lstlisting}
\end{ejercicio}

\begin{solucion}
Insertamos los asertos de Hoare aplicando sucesivamente el axioma de asignación y la regla de composición secuencial:

\begin{lstlisting}[language=C++]
// Estado inicial parametrizado por las constantes c1 y c2:
{ x == c1 && y == c2 }

x = x + y;

{ x == c1 + c2 && y == c2 }

y = x * y;

{ x == c1 + c2 && y == (c1 + c2) * c2 }

x = x - y;

{ x == (c1 + c2) - (c1 + c2) * c2 && y == (c1 + c2) * c2 }
\end{lstlisting}

Factorizando algebraicamente el aserto final:
\begin{align*}
x_{\text{final}} &= (c_1 + c_2) - (c_1 + c_2) \cdot c_2 = \mathbf{(c_1 + c_2)(1 - c_2)}, \\
y_{\text{final}} &= \mathbf{(c_1 + c_2) \cdot c_2}.
\end{align*}
Por tanto, el aserto final que demuestra la corrección total del programa es:
\[
\mathbf{\left\{ x = (c_1 + c_2)(1 - c_2) \;\land\; y = (c_1 + c_2) \cdot c_2 \right\}}.
\]
\end{solucion}

% ------------------------------------------------------------------------------
% EJERCICIO 20
% ------------------------------------------------------------------------------
\begin{ejercicio}[20: Demostración de suma concurrente de potencias de 2]
Demostrar que el siguiente triple es cierto:
\begin{lstlisting}[language=C++]
{ x == 0 }
cobegin
    < x = x + 1 >;
 ||
    < x = x + 2 >;
 ||
    < x = x + 4 >;
coend
{ x == 7 }
\end{lstlisting}
\end{ejercicio}

\begin{solucion}
Denotamos los procesos atómicos como:
\[
S_1 \equiv \langle x := x + 1 \rangle, \qquad S_2 \equiv \langle x := x + 2 \rangle, \qquad S_3 \equiv \langle x := x + 4 \rangle.
\]
Notemos que cada proceso suma una potencia de 2 ($2^0, 2^1, 2^2$). El valor de $x$ codifica directamente en binario qué procesos han finalizado:
\begin{itemize}
    \item $P_1$ conmuta el bit 0 ($+1$).
    \item $P_2$ conmuta el bit 1 ($+2$).
    \item $P_3$ conmuta el bit 2 ($+4$).
\end{itemize}

\textbf{Precondiciones Locales (condición de que su bit respectivo aún vale 0):}
\begin{align*}
\text{pre}(S_1) &\equiv \{x \in \{0, 2, 4, 6\}\}, \\
\text{pre}(S_2) &\equiv \{x \in \{0, 1, 4, 5\}\}, \\
\text{pre}(S_3) &\equiv \{x \in \{0, 1, 2, 3\}\}.
\end{align*}

\textbf{Poscondiciones Locales (su bit respectivo ha pasado a valer 1):}
\begin{align*}
\text{post}(S_1) &\equiv \{x \in \{1, 3, 5, 7\}\}, \\
\text{post}(S_2) &\equiv \{x \in \{2, 3, 6, 7\}\}, \\
\text{post}(S_3) &\equiv \{x \in \{4, 5, 6, 7\}\}.
\end{align*}

\textbf{Comprobación de No Interferencia de Owicki-Gries:}
Debemos comprobar que la ejecución de $S_2$ o $S_3$ no invalida ni la precondición ni la poscondición de $S_1$.
Tomemos por ejemplo $\text{pre}(S_1) \equiv x \in \{0, 2, 4, 6\}$. Si se ejecuta $S_2$ ($x := x + 2$):
\begin{itemize}
    \item Si $x = 0 \to x = 2 \in \{0, 2, 4, 6\}$.
    \item Si $x = 4 \to x = 6 \in \{0, 2, 4, 6\}$.
\end{itemize}
El aserto $x \in \{0, 2, 4, 6\}$ se preserva estrictamente invariantemente bajo $S_2$. El mismo razonamiento se aplica simétricamente a todas las demás combinaciones de procesos y asertos.

\textbf{Poscondición Global Concurrente:}
Al concluir el \texttt{coend}, se verifica la intersección de todas las poscondiciones locales:
\[
x \in \text{post}(S_1) \cap \text{post}(S_2) \cap \text{post}(S_3) = \{1, 3, 5, 7\} \cap \{2, 3, 6, 7\} \cap \{4, 5, 6, 7\} = \mathbf{\{7\}}.
\]
Queda demostrado formalmente que el triple concurrente concluye con $\{x = 7\}$.
\end{solucion}

% ------------------------------------------------------------------------------
% EJERCICIO 21
% ------------------------------------------------------------------------------
\begin{ejercicio}[21: Demostración de invariancia con cobegin]
Dada la siguiente construcción de composición concurrente $P$:
\begin{lstlisting}[language=C++]
cobegin
    < x = x - 1 >; < x = x + 1 >;
 ||
    < y = y + 1 >; < y = y - 1 >;
coend;
\end{lstlisting}
demostrar que se cumple la invariancia de $\{x == y\}$, es decir, que $\{x == y\} \; P \; \{x == y\}$ es un triple cierto.
\end{ejercicio}

\begin{solucion}
\textbf{1. Corrección Secuencial de cada Rama:}
Para cada proceso por separado, verificamos la preservación de la igualdad $x = y$:
\begin{itemize}
    \item \textbf{Rama 1:}
    \[
    \{x = y\} \; \langle x := x - 1 \rangle \; \{x + 1 = y\} \; \langle x := x + 1 \rangle \; \{x = y\}.
    \]
    Por el axioma de asignación: $(x - 1) + 1 = y \iff x = y$. Demostrado.
    \item \textbf{Rama 2:}
    \[
    \{x = y\} \; \langle y := y + 1 \rangle \; \{x = y - 1\} \; \langle y := y - 1 \rangle \; \{x = y\}.
    \]
    Por el axioma de asignación: $x = (y + 1) - 1 \iff x = y$. Demostrado.
\end{itemize}

\textbf{2. Demostración de No Interferencia de Owicki-Gries:}
El conjunto de variables escritas por el proceso 1 es $E(P_1) = \{x\}$, mientras que el proceso 2 escribe en $E(P_2) = \{y\}$.
Para verificar que los pasos intermedios no interfieren destructivamente:
\begin{itemize}
    \item Si se intercala un paso de $P_2$ cuando $P_1$ está en su estado intermedio $\{x + 1 = y\}$, $P_2$ ejecuta $\langle y := y + 1 \rangle$, llevando el estado a $\{x + 2 = y\}$.
    \item Como ambos procesos realizan operaciones complementarias ($+1$ y $-1$) en sus variables privadas respectivas, la diferencia neta $\Delta x - \Delta y$ al finalizar ambos es estrictamente $(0) - (0) = 0$.
\end{itemize}
Formalmente, aplicando variables auxiliares o el invariante global $I \equiv (x_{\text{inicial}} - y_{\text{inicial}} = x - y + \delta_1 - \delta_2)$, al finalizar ambos procesos se tiene $\delta_1 = 0$ y $\delta_2 = 0$, concluyendo:
\[
x - y = x_{\text{inicial}} - y_{\text{inicial}} = 0 \implies \mathbf{\{x = y\}}.
\]
\end{solucion}

% ==============================================================================
% BLOQUE V: REGLAS DE DERIVACIÓN DE HOARE, CONDICIONALES Y ALIAS
% ==============================================================================
\section{Bloque V: Reglas de Derivación de Hoare, Condicionales y Alias}

% ------------------------------------------------------------------------------
% EJERCICIO 22
% ------------------------------------------------------------------------------
\begin{ejercicio}[22: Regla de la conjunción]
Usando la regla de la conjunción de la lógica de Hoare, demostrar que:
\[
\{i > 2\} \; i := 2 * i \; \{i > 4\}
\]
\end{ejercicio}

\begin{solucion}
La \textbf{Regla de la Conjunción} establece que si un programa $S$ satisface dos especificaciones con la misma precondición, satisface simultáneamente su conjunción:
\[
\frac{\{P\} \; S \; \{Q\}, \qquad \{P\} \; S \; \{R\}}{\{P\} \; S \; \{Q \land R\}}
\]

Para demostrar el triple requerido, consideramos los dos triples siguientes:
\begin{enumerate}
    \item \textbf{Triple 1:} $\{i > 2\} \; i := 2 * i \; \{i > 2\}$
    Aplicando el axioma de asignación:
    \[
    \{i > 2\}_{2i}^i \equiv \{2i > 2\} \equiv \{i > 1\}.
    \]
    Como $\{i > 2\} \implies \{i > 1\}$, por fortalecimiento de precondición el triple es cierto.
    
    \item \textbf{Triple 2:} $\{V\} \; i := 2 * i \; \{i \text{ es par}\}$. O bien, con la variable inicial $i_0$:
    \[
    \{i = i_0 \land i_0 > 2\} \; i := 2 * i \; \{i = 2 \cdot i_0\}.
    \]
    Directamente por el axioma de asignación: $\{2i = 2i_0\} \equiv \{i = i_0\}$.
\end{enumerate}

Aplicando la regla de la conjunción entre $\{i > 2\} \; i := 2 * i \; \{i = 2 \cdot i_0\}$ y la condición $i_0 > 2$:
\[
\{i = i_0 \land i_0 > 2\} \; i := 2 * i \; \{i = 2 \cdot i_0 \land i_0 > 2\}.
\]
Dado que $(i = 2 \cdot i_0 \land i_0 > 2) \implies (i > 4)$, por la regla de debilitamiento de poscondición concluimos:
\[
\mathbf{\{i > 2\} \; i := 2 * i \; \{i > 4\}}.
\]
\end{solucion}

% ------------------------------------------------------------------------------
% EJERCICIO 23
% ------------------------------------------------------------------------------
\begin{ejercicio}[23: Intercambio y combinación de variables sin variable auxiliar]
Sean $A$ y $B$ los valores iniciales de $a$ y $b$ respectivamente. Escribir un fragmento de código que tenga $\{a = A + B \land b = A - B\}$ como poscondición y demostrar formalmente que el código es correcto.
\end{ejercicio}

\begin{solucion}
Queremos diseñar un bloque de sentencias $C$ tal que:
\[
\{a = A \land b = B\} \; C \; \{a = A + B \land b = A - B\}
\]
sin recurrir a ninguna variable temporal auxiliar.

\textbf{1. Propuesta de Código $C$:}
\begin{lstlisting}[language=C++]
a = a + b;
b = a - 2*b;
\end{lstlisting}

\textbf{2. Demostración Formal Paso a Paso (Sustitución hacia atrás):}
Aplicamos la regla de composición secuencial:

\begin{itemize}[leftmargin=*]
    \item \textbf{Paso 2 (Segunda instrucción $b := a - 2*b$):}
    Queremos que la poscondición sea $Q \equiv \{a = A + B \land b = A - B\}$.
    Por el axioma de asignación, sustituimos $b$ por $(a - 2b)$ en $Q$:
    \begin{align*}
    Q[a - 2b / b] &\equiv \{a = A + B \;\land\; a - 2b = A - B\} \\
    &\equiv \{a = A + B \;\land\; (A + B) - 2b = A - B\} \\
    &\equiv \{a = A + B \;\land\; 2B = 2b\} \\
    &\equiv \{a = A + B \;\land\; b = B\}.
    \end{align*}
    Por tanto, el triple para la segunda instrucción es:
    \[
    \{a = A + B \land b = B\} \; b := a - 2*b \; \{a = A + B \land b = A - B\}.
    \]

    \item \textbf{Paso 1 (Primera instrucción $a := a + b$):}
    Tomamos como poscondición intermedia $R \equiv \{a = A + B \land b = B\}$.
    Por el axioma de asignación, sustituimos $a$ por $(a + b)$ en $R$:
    \begin{align*}
    R[a + b / a] &\equiv \{a + b = A + B \;\land\; b = B\} \\
    &\equiv \{a + B = A + B \;\land\; b = B\} \\
    &\equiv \{a = A \;\land\; b = B\}.
    \end{align*}
    Por tanto, el triple para la primera instrucción es:
    \[
    \{a = A \land b = B\} \; a := a + b \; \{a = A + B \land b = B\}.
    \]
\end{itemize}

Por la regla de composición secuencial, uniendo ambos pasos queda demostrado que el código es estrictamente correcto.
\end{solucion}

% ------------------------------------------------------------------------------
% EJERCICIO 24
% ------------------------------------------------------------------------------
\begin{ejercicio}[24: Demostración formal de sentencia condicional]
Demostrar que el siguiente triple de Hoare es correcto:
\[
\{V\} \quad \text{if } a > 0 \text{ then } x := a \text{ else } x := -a \quad \{x \ge 0 \land x^2 = a^2\}
\]
\end{ejercicio}

\begin{solucion}
La \textbf{Regla del Condicional} (sentencia \texttt{if-then-else}) establece:
\[
\frac{\{P \land B\} \; S_1 \; \{Q\}, \qquad \{P \land \neg B\} \; S_2 \; \{Q\}}{\{P\} \; \text{if } B \text{ then } S_1 \text{ else } S_2 \; \{Q\}}
\]
En este problema: $P \equiv V$, $B \equiv (a > 0)$, $S_1 \equiv (x := a)$, $S_2 \equiv (x := -a)$, y la poscondición objetivo es $Q \equiv \{x \ge 0 \land x^2 = a^2\}$.

\begin{enumerate}[leftmargin=*]
    \item \textbf{Demostración de la rama THEN ($S_1$):}
    Debemos probar $\{V \land a > 0\} \; x := a \; \{x \ge 0 \land x^2 = a^2\}$.
    Aplicamos el axioma de asignación sobre $Q$:
    \[
    Q[a / x] \equiv \{a \ge 0 \land a^2 = a^2\} \equiv \{a \ge 0\}.
    \]
    Como la precondición de la rama es $\{a > 0\}$, y trivialmente $(a > 0) \implies (a \ge 0)$, por la regla de fortalecimiento de precondición queda probado el primer triple:
    \[
    \{a > 0\} \; x := a \; \{x \ge 0 \land x^2 = a^2\}.
    \]

    \item \textbf{Demostración de la rama ELSE ($S_2$):}
    Debemos probar $\{V \land \neg(a > 0)\} \equiv \{a \le 0\} \; x := -a \; \{x \ge 0 \land x^2 = a^2\}$.
    Aplicamos el axioma de asignación sobre $Q$:
    \[
    Q[-a / x] \equiv \{-a \ge 0 \land (-a)^2 = a^2\} \equiv \{a \le 0 \land a^2 = a^2\} \equiv \{a \le 0\}.
    \]
    La precondición generada por el axioma coincide exactamente con $\{a \le 0\}$, por lo que el triple es un axioma directo.
\end{enumerate}

Habiéndose probado ambas premisas, la regla del condicional concluye directamente la corrección del programa.
\end{solucion}

% ------------------------------------------------------------------------------
% EJERCICIO 25
% ------------------------------------------------------------------------------
\begin{ejercicio}[25: Demostración de secuencia de asignaciones]
Demostrar formalmente utilizando los axiomas de la lógica de Hoare que:
\[
\{i \cdot j + 2j + 3i = 0\} \quad j := j + 3; \; i := i + 2 \quad \{i \cdot j = 6\}
\]
\end{ejercicio}

\begin{solucion}
Sea el programa compuesto $S \equiv S_1; S_2$ con $S_1 \equiv (j := j + 3)$ y $S_2 \equiv (i := i + 2)$, y la poscondición final $Q \equiv \{i \cdot j = 6\}$.

\textbf{Paso 1: Cálculo del aserto intermedio para la segunda instrucción ($S_2$):}
Por el axioma de asignación:
\[
Q[i + 2 / i] \equiv \{(i + 2) \cdot j = 6\} \equiv \{i \cdot j + 2j = 6\}.
\]
Por tanto, tenemos el triple demostrado:
\[
\{i \cdot j + 2j = 6\} \; i := i + 2 \; \{i \cdot j = 6\}.
\]

\textbf{Paso 2: Cálculo de la precondición inicial para la primera instrucción ($S_1$):}
Tomamos como poscondición requerida el aserto intermedio $R \equiv \{i \cdot j + 2j = 6\}$.
Aplicamos el axioma de asignación sustituyendo $j$ por $(j + 3)$ en $R$:
\begin{align*}
R[j + 3 / j] &\equiv \{i \cdot (j + 3) + 2 \cdot (j + 3) = 6\} \\
&\equiv \{i \cdot j + 3i + 2j + 6 = 6\} \\
&\equiv \{i \cdot j + 2j + 3i = 0\}.
\end{align*}
Por tanto, tenemos el triple demostrado:
\[
\{i \cdot j + 2j + 3i = 0\} \; j := j + 3 \; \{i \cdot j + 2j = 6\}.
\]

\textbf{Paso 3: Aplicación de la regla de composición:}
Uniendo ambos triples mediante la regla de composición secuencial:
\[
\frac{\{i \cdot j + 2j + 3i = 0\} \; j := j + 3 \; \{i \cdot j + 2j = 6\}, \quad \{i \cdot j + 2j = 6\} \; i := i + 2 \; \{i \cdot j = 6\}}{\{i \cdot j + 2j + 3i = 0\} \quad j := j + 3; \; i := i + 2 \quad \{i \cdot j = 6\}}
\]
El resultado coincide exactamente con la precondición dada en el enunciado. Q.E.D.
\end{solucion}

% ------------------------------------------------------------------------------
% EJERCICIO 26
% ------------------------------------------------------------------------------
\begin{ejercicio}[26: Guarda booleana en la regla del bucle while]
¿Por qué en la regla del bucle \texttt{while B do}, la condición $B$ debe ser verdadera al comienzo de cada ejecución del cuerpo del bucle?
\end{ejercicio}

\begin{solucion}
La regla de inferencia de Hoare para el bucle \texttt{while} (\textit{regla de iteración}) se formula como:
\[
\frac{\{I \land B\} \; S \; \{I\}}{\{I\} \; \text{while } B \text{ do } S \; \{I \land \neg B\}}
\]
donde $I$ representa el invariante de bucle, $B$ es la guarda o condición booleana y $S$ es el cuerpo del bucle.

\textbf{Razones fundamentales:}
\begin{enumerate}
    \item \textbf{Semántica operacional del bucle:} Por definición, el cuerpo $S$ solo llega a ejecutarse si y solo si la evaluación de la guarda booleana $B$ resulta \texttt{true}. Si $B$ fuese falsa, el flujo de ejecución omite inmediatamente el cuerpo $S$ y transfiere el control a la instrucción posterior.
    \item \textbf{Precisión del invariante inductivo:} Al comenzar cualquier iteración de $S$, sabemos con total certeza que el estado actual satisface no solo el invariante general $I$, sino también la condición de entrada $B$. Esta premisa adicional $\{I \land B\}$ es estrictamente necesaria para poder demostrar que las transformaciones realizadas por $S$ logran restablecer el invariante $I$ al final de la iteración.
    \item \textbf{Garantía de salida:} Al terminar el bucle, la condición de parada garantiza que la última evaluación de la guarda fue falsa ($\neg B$). Puesto que cada iteración preservó $I$, el aserto resultante final es forzosamente la conjunción $\{I \land \neg B\}$.
\end{enumerate}
\end{solucion}

% ------------------------------------------------------------------------------
% EJERCICIO 27
% ------------------------------------------------------------------------------
\begin{ejercicio}[27: El problema de los alias en la lógica de programas]
Considerar una función con dos argumentos que se usa en un programa. Explicar por qué el uso de \textbf{alias} (\textit{aliasing}) puede ser un problema grave para la verificación formal con la Lógica de Hoare.
\end{ejercicio}

\begin{solucion}
El fenómeno de \textbf{aliasing} (alias) ocurre cuando dos o más identificadores, nombres de variables o punteros/referencias diferentes hacen referencia exactamente a la misma posición física de memoria en tiempo de ejecución.

\textbf{Por qué rompe la Lógica de Hoare:}
El axioma fundamental de asignación de Hoare:
\[
\{P[e/x]\} \; x := e \; \{P\}
\]
se basa en la hipótesis implícita de \textbf{independencia sintáctica}: asume que modificar la variable $x$ \textbf{no afecta a ninguna otra variable} presente en la aserción $P$.

\textbf{Ejemplo Demostrativo:}
Supongamos un procedimiento con dos argumentos pasados por referencia:
\begin{lstlisting}[language=C++]
procedure duplicar(var x : integer; var y : integer)
begin
  x := 10;
  y := 20;
end;
\end{lstlisting}
Si suponemos variables distintas, la lógica estándar deduciría:
\[
\{V\} \quad x := 10; \; y := 20 \quad \{x = 10 \land y = 20\}.
\]
Sin embargo, si el invocador llama a \texttt{duplicar(a, a)}, los parámetros $x$ e $y$ son \textbf{alias} de la misma variable $a$.
\begin{itemize}
    \item La sentencia $y := 20$ sobreescribe silenciosamente la posición referenciada por $x$.
    \item El resultado real en memoria es $a = 20$, lo que hace que la poscondición formal $\{x = 10 \land y = 20\}$ sea \textbf{FALSA} ($20 = 10 \land 20 = 20 \implies \text{Falso}$).
\end{itemize}
Por esta razón, la semántica axiomática clásica prohíbe el aliasing o exige modelar la memoria de forma explícita mediante funciones de estado de montón (\textit{heap/store semantics}) o Lógica de Separación (\textit{Separation Logic}).
\end{solucion}

% ==============================================================================
% BLOQUE VI: ASIGNACIÓN A ARRAYS Y PRECONDICIÓN MÁS DÉBIL
% ==============================================================================
\section{Bloque VI: Asignación a Arrays y Precondición Más Débil}

% ------------------------------------------------------------------------------
% EJERCICIO 28
% ------------------------------------------------------------------------------
\begin{ejercicio}[28: Asignación a elementos de un vector]
Demostrar la corrección del siguiente triple:
\[
\{a[i] \ge 0\} \quad a[i] := a[i] + a[j] \quad \{a[i] \ge a[j]\}
\]
\end{ejercicio}

\begin{solucion}
Para razonar sobre elementos de un array, la lógica de Hoare modela la asignación a un elemento de un vector $a[i] := e$ como una actualización de la función que representa el array completo:
\[
a := (a; \; i: e)
\]
donde el array modificado $(a; \; i: e)$ tiene la propiedad:
\[
(a; \; i: e)[k] = \begin{cases}
e & \text{si } k = i \\
a[k] & \text{si } k \neq i
\end{cases}
\]

Queremos demostrar $\{a[i] \ge 0\} \; a[i] := a[i] + a[j] \; \{a[i] \ge a[j]\}$.
Aplicando la sustitución formal de arrays en la poscondición $Q \equiv \{a[i] \ge a[j]\}$:
\[
Q[(a; \; i: a[i] + a[j]) / a] \equiv \left\{ (a; \; i: a[i] + a[j])[i] \;\ge\; (a; \; i: a[i] + a[j])[j] \right\}.
\]
Por definición, el término de la izquierda siempre es $(a; \; i: a[i] + a[j])[i] = a[i] + a[j]$.
Para el término de la derecha, distinguimos dos casos según si los índices coinciden o no:

\begin{itemize}[leftmargin=*]
    \item \textbf{Caso 1 ($i = j$, mismo elemento):}
    En este caso, $(a; \; i: a[i] + a[j])[j] = (a; \; i: a[i] + a[i])[i] = a[i] + a[i]$.
    La condición requerida resulta:
    \[
    a[i] + a[i] \ge a[i] + a[i] \iff 0 \ge 0 \quad (\text{Tautología, siempre verdadero } V).
    \]
    Como $\{a[i] \ge 0\} \implies V$, se cumple trivialmente.

    \item \textbf{Caso 2 ($i \neq j$, elementos distintos):}
    Dado que $j \neq i$, el elemento en la posición $j$ no sufre modificación alguna:
    \[
    (a; \; i: a[i] + a[j])[j] = a[j].
    \]
    La condición requerida resulta:
    \[
    a[i] + a[j] \ge a[j] \iff a[i] \ge 0.
    \]
    Esta expresión coincide exactamente con la precondición dada en el enunciado $\{a[i] \ge 0\}$.
\end{itemize}
En ambos casos, la precondición $\{a[i] \ge 0\}$ garantiza que tras la asignación se verificará $a[i] \ge a[j]$. Queda demostrado.
\end{solucion}

% ------------------------------------------------------------------------------
% EJERCICIO 29
% ------------------------------------------------------------------------------
\begin{ejercicio}[29: Cálculo de Precondición más débil con arrays]
Hallar la precondición $\{P\}$ que hace que el siguiente triple sea correcto:
\[
\{P\} \quad a[i] := 2 * b \quad \{j \le i \;\land\; k < i \;\land\; a[i] + a[j - 1] + a[k] > b\}
\]
\end{ejercicio}

\begin{solucion}
La instrucción asigna al elemento $a[i]$ el valor $2b$, transformando el array $a$ en $(a; \; i: 2b)$.
Por el axioma de asignación para vectores, la precondición más débil se obtiene sustituyendo el array $a$ por $(a; \; i: 2b)$ en la poscondición:
\[
P \equiv \left( j \le i \;\land\; k < i \;\land\; (a; \; i: 2b)[i] + (a; \; i: 2b)[j - 1] + (a; \; i: 2b)[k] > b \right).
\]

Evaluamos cada acceso al array en el nuevo estado:
\begin{enumerate}
    \item $(a; \; i: 2b)[i] = 2b$ (acceso directo a la posición recién escrita).
    \item Para el índice $k$: la poscondición impone $k < i$, por lo que $k \neq i$. Por la propiedad del array actualizado:
    \[
    (a; \; i: 2b)[k] = a[k].
    \]
    \item Para el índice $j - 1$: la poscondición impone $j \le i$, de lo que se deduce que $j - 1 < i$, por lo que $(j - 1) \neq i$. Por consiguiente:
    \[
    (a; \; i: 2b)[j - 1] = a[j - 1].
    \]
\end{enumerate}

Sustituyendo estos valores simplificados en la desigualdad:
\[
2b + a[j - 1] + a[k] > b \iff a[j - 1] + a[k] > b - 2b \iff a[j - 1] + a[k] > -b.
\]

Por tanto, la precondición más débil exacta $\{P\}$ es:
\[
\mathbf{\{P\} \equiv \left\{ j \le i \;\land\; k < i \;\land\; a[j - 1] + a[k] > -b \right\}}.
\]
\end{solucion}

% ==============================================================================
% BLOQUE VII: BUCLES, INVARIANTES Y TERMINACIÓN
% ==============================================================================
\section{Bloque VII: Bucles, Invariantes y Terminación}

% ------------------------------------------------------------------------------
% EJERCICIO 30
% ------------------------------------------------------------------------------
\begin{ejercicio}[30: Demostración de terminación de bucle para n > 0]
Demostrar que para $n > 0$ el siguiente fragmento de programa termina:
\begin{lstlisting}[language=C++]
i := 1;
f := 1;
while (i <> n) do begin
  i := i + 1;
  f := f * i;
end;
\end{lstlisting}
\end{ejercicio}

\begin{solucion}
Para demostrar formalmente la \textbf{terminación} de un bucle se utiliza el Teorema de Terminación basado en una \textbf{función variante} o \textbf{función de cota} $V : \text{Estado} \to \mathbb{Z}$ sobre un conjunto bien fundado (generalmente $\mathbb{N}$).
Debemos probar:
\begin{enumerate}
    \item $V$ está acotada inferiormente mientras el bucle continúe iterando: $\{I \land B\} \implies V \ge 0$.
    \item Cada ejecución del cuerpo del bucle decrementa estrictamente el valor de $V$: si el valor antes de iterar es $V_0$, al final de la iteración se cumple $V < V_0$.
\end{enumerate}

\textbf{1. Definición de la Función de Cota $V$:}
La guarda del bucle es $B \equiv (i \neq n)$. Como $i$ parte de 1 y se incrementa de 1 en 1 hacia $n$, definimos:
\[
V(i) = n - i.
\]

\textbf{2. Demostración de Cota Inferior:}
El invariante del bucle es $I \equiv (1 \le i \le n \land f = i!)$.
Mientras la condición del bucle $B \equiv (i \neq n)$ sea verdadera:
\[
(1 \le i \le n \;\land\; i \neq n) \implies 1 \le i \le n - 1 \implies n - i \ge 1 > 0.
\]
Por tanto, $V(i) = n - i \ge 1 \ge 0$.

\textbf{3. Decrecimiento Estricto en cada Iteración:}
Sea $V_0 = n - i$ el valor de la función de cota al inicio de la iteración.
El cuerpo del bucle ejecuta $i := i + 1$. El nuevo valor de la cota $V'$ es:
\[
V' = n - (i + 1) = (n - i) - 1 = V_0 - 1.
\]
Como $V_0 - 1 < V_0$, la función de cota disminuye estrictamente en 1 en cada iteración.

\textbf{Conclusión:} Como $V(i)$ toma valores en los números enteros, está acotada inferiormente por 0 y decrece estrictamente en cada ciclo, el bucle no puede iterar indefinidamente y termina en exactamente $n - 1$ iteraciones.
\end{solucion}

% ------------------------------------------------------------------------------
% EJERCICIO 31
% ------------------------------------------------------------------------------
\begin{ejercicio}[31: Obtención de poscondiciones apropiadas]
Para cada uno de los siguientes fragmentos de código, obtener la poscondición más apropiada:
\begin{enumerate}[label=(\alph*),leftmargin=*]
    \item $\{i < 10\} \; i := 2 * i + 1;$
    \item $\{i > 0\} \; i := i - 1;$
    \item $\{i > j\} \; i := i + 1; \; j := j + 1;$
    \item $\{V\} \; i := 3; \; j := 2 * i;$
\end{enumerate}
\end{ejercicio}

\begin{solucion}
Aplicando directamente el axioma de asignación y composición secuencial:
\begin{enumerate}[label=(\alph*),leftmargin=*]
    \item \textbf{Fragmento (a):} Si $i_0 < 10$, y la nueva variable es $i = 2i_0 + 1$, despejando $i_0 = (i - 1)/2 < 10 \iff i < 21$.
    Poscondición: $\mathbf{\{i < 21\}}$.
    \item \textbf{Fragmento (b):} Si $i_0 > 0$, con $i_0 \in \mathbb{Z} \implies i_0 \ge 1$. Tras $i = i_0 - 1 \implies i_0 = i + 1 \ge 1 \iff i \ge 0$.
    Poscondición: $\mathbf{\{i \ge 0\}}$ (o equivalentemente $\{i > -1\}$).
    \item \textbf{Fragmento (c):} Ambos operandos se incrementan simultáneamente en 1 unidad:
    \[
    \{i > j\} \; i := i + 1 \; \{i > j + 1\} \; j := j + 1 \; \{i > j\}.
    \]
    Poscondición: $\mathbf{\{i > j\}}$.
    \item \textbf{Fragmento (d):} Con precondición trivial $\{V\}$, tras $i := 3$ se tiene $i = 3$. Luego $j := 2 * 3 = 6$.
    Poscondición: $\mathbf{\{i = 3 \land j = 6\}}$.
\end{enumerate}
\end{solucion}

% ------------------------------------------------------------------------------
% EJERCICIO 32
% ------------------------------------------------------------------------------
\begin{ejercicio}[32: Verificación de código secuencial con reglas de Hoare]
Verificar formalmente el siguiente código indicando todas las reglas y axiomas utilizados:
\[
\{y > 0\} \quad x := x + y; \; x := x - y \quad \{x = x_0 \land y > 0\}
\]
\end{ejercicio}

\begin{solucion}
Sea $x_0, y_0$ los valores iniciales de las variables antes de ejecutar el bloque, con precondición $P \equiv \{y = y_0 \land y_0 > 0 \land x = x_0\}$.
La poscondición buscada es $Q \equiv \{x = x_0 \land y = y_0 \land y > 0\}$.

\textbf{1. Segunda Asignación ($x := x - y$):}
Por el \textbf{Axioma de Asignación}:
\[
Q[x - y / x] \equiv \{x - y = x_0 \land y = y_0 \land y > 0\} \equiv \{x = x_0 + y \land y = y_0 \land y > 0\}.
\]
Denotamos este aserto intermedio como $R$. El triple:
\[
\{x = x_0 + y \land y = y_0 \land y > 0\} \; x := x - y \; \{x = x_0 \land y = y_0 \land y > 0\}
\]
es un axioma.

\textbf{2. Primera Asignación ($x := x + y$):}
Aplicamos el \textbf{Axioma de Asignación} sobre el aserto intermedio $R$:
\[
R[x + y / x] \equiv \{(x + y) = x_0 + y \land y = y_0 \land y > 0\} \equiv \{x = x_0 \land y = y_0 \land y > 0\}.
\]
El triple:
\[
\{x = x_0 \land y = y_0 \land y > 0\} \; x := x + y \; \{x = x_0 + y \land y = y_0 \land y > 0\}
\]
es un axioma.

\textbf{3. Regla de Composición Secuencial:}
Por la regla de composición entre (2) y (1):
\[
\frac{\{P\} \; x := x + y \; \{R\}, \qquad \{R\} \; x := x - y \; \{Q\}}{\{P\} \quad x := x + y; \; x := x - y \quad \{Q\}}
\]
Queda verificado formalmente.
\end{solucion}

% ------------------------------------------------------------------------------
% EJERCICIO 33
% ------------------------------------------------------------------------------
\begin{ejercicio}[33: Verificación de estructura condicional if-then]
Verificar el siguiente código indicando todas las reglas usadas:
\[
\{V\} \quad \text{if } x < 0 \text{ then } x := -x \quad \{x \ge 0\}
\]
\end{ejercicio}

\begin{solucion}
La sentencia \texttt{if B then S} equivale sintácticamente a \texttt{if B then S else null}.
La \textbf{Regla del Condicional} requiere probar dos triples:
\[
\frac{\{V \land x < 0\} \; x := -x \; \{x \ge 0\}, \qquad \{V \land \neg(x < 0)\} \; \text{null} \; \{x \ge 0\}}{\{V\} \; \text{if } x < 0 \text{ then } x := -x \; \{x \ge 0\}}
\]

\textbf{1. Rama THEN ($x < 0$):}
Aplicamos el axioma de asignación sobre la poscondición $Q \equiv \{x \ge 0\}$:
\[
Q[-x / x] \equiv \{-x \ge 0\} \equiv \{x \le 0\}.
\]
Por el axioma de asignación: $\{x \le 0\} \; x := -x \; \{x \ge 0\}$.
Como la precondición de la rama es $\{x < 0\}$ y $(x < 0) \implies (x \le 0)$, por la \textbf{Regla de Fortalecimiento de Precondición}:
\[
\{x < 0\} \; x := -x \; \{x \ge 0\}.
\]

\textbf{2. Rama ELSE ($\neg(x < 0) \equiv x \ge 0$):}
Por el \textbf{Axioma de la Sentencia Vacía} (\texttt{null}):
\[
\{x \ge 0\} \; \text{null} \; \{x \ge 0\}.
\]

Aplicando la regla del condicional a ambas ramas, queda demostrado el triple con poscondición $\{x \ge 0\}$.
\end{solucion}

% ------------------------------------------------------------------------------
% EJERCICIO 34
% ------------------------------------------------------------------------------
\begin{ejercicio}[34: Verificación de un algoritmo de búsqueda con invariante]
Verificar el siguiente segmento de programa para encontrar el máximo de un array $a[1..n]$, suponiendo el invariante dado:
\[
I \equiv \left( \bigwedge_{k=1}^i (\max \ge a[k]) \;\land\; 1 \le i \le n + 1 \right)
\]

\begin{lstlisting}[language=C++]
max := a[1];
i := 1;
while (i <> n + 1) do begin
  if (a[i] >= max) then
    max := a[i];
  i := i + 1;
end;
\end{lstlisting}
poscondición final a demostrar: $\displaystyle \left\{ \bigwedge_{k=1}^n (\max \ge a[k]) \right\}$.
\end{ejercicio}

\begin{solucion}
Para demostrar la corrección parcial del bucle mediante la \textbf{Regla de Iteración de Hoare}:
\[
\frac{\{I \land B\} \; S \; \{I\}}{\{I\} \; \text{while } B \text{ do } S \; \{I \land \neg B\}}
\]
debemos demostrar tres etapas formales:

\textbf{Etapa 1: Inicialización (Establecimiento del invariante antes del bucle):}
Debemos probar $\{n \ge 1\} \; \max := a[1]; \; i := 1 \; \{I\}$.
Por sustitución hacia atrás:
\[
I[1 / i, a[1] / \max] \equiv \left( \bigwedge_{k=1}^1 (a[1] \ge a[k]) \;\land\; 1 \le 1 \le n + 1 \right) \equiv (a[1] \ge a[1] \land 1 \le n + 1) \equiv V.
\]
Por tanto, la inicialización establece válidamente el invariante $I$.

\textbf{Etapa 2: Mantenimiento del Invariante por el cuerpo del bucle:}
La guarda es $B \equiv (i \neq n + 1)$. El estado antes de iterar cumple:
\[
\{I \land B\} \equiv \left( \bigwedge_{k=1}^{i-1} (\max \ge a[k]) \;\land\; \max \ge a[i] \;\text{o no} \;\land\; 1 \le i \le n \right).
\]
Analizamos el condicional:
\begin{itemize}
    \item Si $a[i] \ge \max$, se ejecuta $\max := a[i]$. En el nuevo estado, $\max$ es mayor o igual que el nuevo elemento y, por transitividad, mayor o igual que todos los anteriores.
    \item Si $a[i] < \max$, se omite la asignación. El máximo previo ya era estrictamente mayor que $a[i]$ y mayor o igual que todos los anteriores.
\end{itemize}
En ambos casos, tras la sentencia condicional se satisface:
\[
\bigwedge_{k=1}^i (\max \ge a[k]).
\]
A continuación, la instrucción $i := i + 1$ actualiza el contador. Al sustituir $i$ por $i + 1$, el rango cubierto es exactamente $\{1, \dots, (i' - 1)\}$, preservando la forma del invariante $I$ para la siguiente iteración.

\textbf{Etapa 3: Poscondición al finalizar el bucle:}
Al terminar el bucle se cumple $\{I \land \neg B\}$:
\[
\neg(i \neq n + 1) \iff i = n + 1.
\]
Sustituyendo $i = n + 1$ en el invariante $I$:
\[
I \land (i = n + 1) \implies \bigwedge_{k=1}^{n+1 - 1} (\max \ge a[k]) \equiv \mathbf{\bigwedge_{k=1}^n (\max \ge a[k])}.
\]
La poscondición coincide exactamente con la especificación del máximo de los $n$ elementos del array.
\end{solucion}

% ------------------------------------------------------------------------------
% EJERCICIO 35
% ------------------------------------------------------------------------------
\begin{ejercicio}[35: Evaluación del número combinatorio con bucle único]
Dados $n \ge 0, i \le n$, demostrar que el siguiente segmento de programa evalúa el coeficiente binomial $\binom{n}{i} = \frac{n!}{i! (n - i)!}$ dado el invariante:
\begin{align*}
I \equiv \big( &(i > k \;\lor\; afact = i!) \;\land\; (n - i > k \;\lor\; bfact = (n - i)!) \;\land \\
& fact = k! \;\land\; 0 \le k \le n \big)
\end{align*}

\begin{lstlisting}[language=C++]
k := 0; fact := 1;
while (k <> n) do begin
  k := k + 1;
  fact := fact * k;
  if (k <= i) then afact := fact;
  if (k <= n - i) then bfact := fact;
end;
bcof := fact / (afact * bfact);
\end{lstlisting}
\end{ejercicio}

\begin{solucion}
\textbf{1. Verificación de la Inicialización:}
Antes de entrar al bucle: $k = 0, fact = 1 = 0!$.
Como $i \ge 0$ y $n - i \ge 0$:
\begin{itemize}
    \item Si $i > 0$, la disyunción $(i > 0 \lor afact = i!)$ es inmediatamente verdadera por el primer término ($i > 0$). Si $i = 0$, $afact$ se asignará adecuadamente o se evalúa con el factorial vacío $0! = 1$.
    \item Análogamente para $n - i \ge 0$.
\end{itemize}
El invariante $I$ se satisface formalmente antes de comenzar.

\textbf{2. Mantenimiento del Invariante en el Bucle:}
En cada iteración, $k$ se incrementa en 1 ($k' = k + 1$) y $fact' = fact \times k' = k! \times k' = k'!$.
\begin{itemize}
    \item Mientras $k' \le i$, el condicional actualiza $afact := fact' = k'!$. Cuando $k'$ alcanza exactamente el valor $i$, se ejecuta la asignación que deja fijado $afact = i!$. Para las iteraciones posteriores ($k' > i$), la condición $k' \le i$ es falsa, por lo que $afact$ nunca vuelve a modificarse, preservando eternamente $afact = i!$.
    \item El mismo razonamiento se aplica idénticamente a $bfact$: cuando $k'$ alcanza $n - i$, se asigna $bfact = (n - i)!$, valor que permanece congelado hasta el final.
\end{itemize}
Por tanto, el invariante $I$ se mantiene inductivamente en cada iteración.

\textbf{3. Terminación y Cálculo del Resultado:}
El bucle finaliza cuando $\neg(k \neq n) \iff k = n$.
Sustituyendo $k = n$ en el invariante $I$:
\begin{itemize}
    \item Como $i \le n$, la condición $i > n$ es falsa, obligando a que se cumpla el segundo término de la disyunción: $\mathbf{afact = i!}$.
    \item Como $n - i \le n$, análogamente se cumple: $\mathbf{bfact = (n - i)!}$.
    \item El valor acumulado del factorial completo es: $\mathbf{fact = n!}$.
\end{itemize}

Por último, la asignación final calcula:
\[
bcof := \frac{fact}{afact \times bfact} = \mathbf{\frac{n!}{i! (n - i)!}} = \mathbf{\binom{n}{i}}.
\]
Queda demostrada la corrección total del programa.
\end{solucion}

% ------------------------------------------------------------------------------
% EJERCICIO 36
% ------------------------------------------------------------------------------
\begin{ejercicio}[36: Demostración de terminación del algoritmo de búsqueda de máximo]
Demostrar la terminación del fragmento de programa dado en el problema 34 anterior. ¿Qué condición se debe imponer sobre el parámetro $n$ para poder realizar la demostración formal?
\end{ejercicio}

\begin{solucion}
\textbf{1. Condición necesaria sobre el parámetro $n$:}
Para que el bucle termine de forma determinista, se debe imponer la condición:
\[
\mathbf{n \ge 0} \quad (\text{y en particular } \mathbf{n \ge 1} \text{ para que el acceso inicial } a[1] \text{ esté bien definido}).
\]
\textbf{Justificación del contraejemplo si $n < 0$:}
El bucle inicializa $i := 1$ y su guarda de continuación es $(i \neq n + 1)$. Como el cuerpo incrementa $i$ estrictamente de uno en uno ($i \leftarrow i + 1$):
\begin{itemize}
    \item La secuencia de valores generada por $i$ es $\{1, 2, 3, 4, \dots\}$.
    \item Si $n < 0$, entonces la meta $n + 1 \le 0$. Como $i$ parte de $1 > n + 1$ y crece positivamente, $i$ \textbf{jamás igualará a $n + 1$}, produciéndose un \textbf{bucle infinito} (divergencia).
    \item Por consiguiente, $n \ge 0$ es la precondición mínima indispensable para garantizar convergencia.
\end{itemize}

\textbf{2. Demostración Formal con Función Variante (Cota):}
Bajo la precondición $n \ge 0$, definimos la función variante:
\[
V(i) = (n + 1) - i.
\]

\textbf{Comprobación de las dos condiciones de terminación:}
\begin{enumerate}
    \item \textbf{Cota inferior positiva mientras el bucle ejecuta:}
    Mientras la guarda $(i \neq n + 1)$ sea verdadera, y dado que por el invariante $1 \le i \le n + 1$, se tiene $i \le n$, luego:
    \[
    V(i) = n + 1 - i \ge n + 1 - n = 1 > 0.
    \]
    La función de cota está estrictamente acotada inferiormente por 0 en todos los estados admisibles.
    
    \item \textbf{Decrecimiento monótono estricto:}
    Al inicio de la iteración, el valor de la función de cota es $V_0 = n + 1 - i$.
    La sentencia $i := i + 1$ modifica el valor a:
    \[
    V' = n + 1 - (i + 1) = (n + 1 - i) - 1 = V_0 - 1.
    \]
    Puesto que $V' < V_0$, la función variante decrementa exactamente en 1 en cada paso.
\end{enumerate}

Por el principio de inducción bien fundada, el bucle concluye obligatoriamente tras exactamente $n$ iteraciones.
\end{solucion}
